\documentclass[11pt,a4paper]{article}

\usepackage[T1]{fontenc}
\usepackage[utf8]{inputenc}
\usepackage[a4paper,margin=27mm]{geometry}
\usepackage{amsmath,amssymb,amsthm,mathtools}
\usepackage{newpxtext,newpxmath}
\usepackage{microtype}
\usepackage{booktabs,tabularx,array}
\usepackage{enumitem}
\usepackage{xcolor}
\usepackage{titlesec}
\usepackage[authoryear,round]{natbib}
\usepackage[colorlinks=true,linkcolor=blue!45!black,
  citecolor=red!45!black,urlcolor=blue!45!black]{hyperref}
\usepackage[nameinlink,noabbrev]{cleveref}

\definecolor{otblue}{HTML}{1F4E79}
\titleformat{\section}{\normalfont\Large\bfseries\color{otblue}}
  {\thesection}{.7em}{}
\titleformat{\subsection}{\normalfont\large\bfseries\color{otblue!80!black}}
  {\thesubsection}{.65em}{}
\titleformat{\paragraph}[runin]
  {\normalfont\bfseries\color{otblue!75!black}}{}{0pt}{}
\titlespacing*{\section}{0pt}{2.7ex plus .7ex minus .2ex}{1.1ex}
\titlespacing*{\subsection}{0pt}{2ex plus .5ex minus .2ex}{.75ex}
\titlespacing*{\paragraph}{0pt}{1.4ex plus .3ex minus .1ex}{.55em}

\newtheorem{theorem}{Theorem}[section]
\newtheorem{proposition}[theorem]{Proposition}
\newtheorem{lemma}[theorem]{Lemma}
\newtheorem{corollary}[theorem]{Corollary}
\theoremstyle{definition}
\newtheorem{definition}[theorem]{Definition}
\newtheorem{example}[theorem]{Example}
\theoremstyle{remark}
\newtheorem{remark}[theorem]{Remark}

\newcommand{\R}{\mathbb R}
\newcommand{\Pcal}{\mathcal P}
\newcommand{\Hcal}{\mathcal H}
\newcommand{\Lcal}{\mathcal L}
\newcommand{\Ecal}{\mathcal E}
\newcommand{\Ccal}{\mathcal C}
\newcommand{\Piopt}{\Pi}
\newcommand{\Id}{\operatorname{Id}}
\newcommand{\supp}{\operatorname{supp}}
\newcommand{\rank}{\operatorname{rank}}
\newcommand{\Ran}{\operatorname{Ran}}
\newcommand{\card}{\operatorname{card}}
\newcommand{\HS}{\operatorname{HS}}
\newcommand{\argmin}{\operatorname*{arg\,min}}
\newcommand{\argmax}{\operatorname*{arg\,max}}
\newcommand{\dd}{\,\mathrm d}
\newcommand{\norm}[1]{\left\lVert #1\right\rVert}
\newcommand{\abs}[1]{\left\lvert #1\right\rvert}
\newcommand{\ip}[2]{\left\langle #1,#2\right\rangle}
\newcommand{\GW}{\operatorname{GW}}
\newcommand{\GM}{\operatorname{GM}}

\hypersetup{
  pdftitle={Toward a Brenier Theorem for Power-Distance Gromov--Wasserstein Transport},
  pdfauthor={Gabriel Peyre},
  pdfsubject={Self-consistent costs, quotient lifting, and Monge maps for Gromov--Wasserstein transport},
  pdfkeywords={Gromov--Wasserstein, Monge map, Brenier theorem, quotient transport, negative type, RKHS, twist condition}
}

\title{\textbf{Toward a Brenier Theorem for\\
Power-Distance Gromov--Wasserstein Transport}\\
\large Self-consistent costs, quotient lifting, and finite transport branches}
\author{Gabriel Peyr\'e}
\date{July 2026}

\begin{document}
\maketitle

\begin{abstract}
Let \(0<p\leq2\), and consider Gromov--Wasserstein transport with distortion
\(
 (\norm{x-x'}^p-\norm{y-y'}^p)^2.
\)
Schoenberg's theorem realizes each power distance as a squared distance in a
Hilbert feature space.  After centering these features, the GW objective
depends on a coupling through one scalar mixed moment and one Hilbert--Schmidt
cross-covariance operator.  This yields an exact representation as ordinary
optimal transport with a learned cross-space cost, regularized by the squared
norm of a tensor-product RKHS.  Every GW minimizer therefore solves a linear
OT problem for its self-consistent learned cost.

This reduction reveals two distinct routes to a Brenier-type result.  First,
an \(m\)-twisted self-consistent cost forces every conditional plan to have at
most \(m\) atoms, hence represents a prescribed GW optimizer as a measurable
mixture of at most \(m\) maps.  Second, conditional concavity gives a
replacement principle: every optimizer of the frozen linear problem remains
GW-optimal.  A twisted transport problem can then be solved on a quotient and
lifted through nonatomic source fibers, producing a genuine Monge optimizer
even when the original twist fibers are infinite.  In particular, if the
learned operator has rank \(R\), absolute continuity of its
\((R+1)\)-dimensional quotient and atomlessness along its fibers imply
existence of a GW-optimal map; a coarea criterion makes these hypotheses
checkable.

We also give differential, multijet-transversality, and analytic-definable
criteria for finite twist.  Analyticity alone is shown to be insufficient.
For \(p=2\), the feature spaces are finite dimensional and the two mechanisms
combine into a rank-free conclusion: for compactly supported marginals, if the
source lies in \(\R^n\), the target lies in \(\R^d\), \(n>d\), and the source
has a density, then squared-distance GW admits a Monge minimizer.  More
precisely, a full-row-rank optimizer is the unique optimizer of its frozen OT
problem and has both two-map and map/anti-map structure; when \(n>d\), it is
itself graph-supported.  A lower-rank optimizer can instead be replaced by a
graph coupling with the same cross-covariance.  Density of graph couplings
already makes the directed Gromov--Monge infimum equal to the GW minimum; the
theorem upgrades this value equality to attainment.  The same direct graph
conclusion holds without a strict dimension gap when the target lies on a
centered sphere.  In equal dimensions, the second full-rank preimage is exactly
the reflection of the first across a source-dependent affine hyperplane;
excluding these reflected pairs on the dual contact fibers gives an exact
twist-based map criterion.  At rank \(d-1\), injectivity of the radial-linear
target feature gives a complementary criterion.  Without such geometry, the
learned cost is only two-twisted.  This recovers the mechanism behind the known
optimal two-map theorem.  The equal-dimensional rotationally invariant case
is completely rigid: Sturm's theorem identifies every optimizer as the
monotone radial quantile map followed by an orthogonal transformation.  On
the line, a more flexible marginal-only quantile criterion gives a global
monotone GW optimizer whenever the radial mixed moment and covariance
magnitude are maximized by the same rearrangement.  The note separates proved
statements, known results, and the open step for \(0<p<2\).
\end{abstract}

\section{Problem, scope, and known results}
\label{sec:setup}

The goal is to understand when the relaxation from deterministic
correspondences to couplings is exact \emph{at the level of minimizers}.  The
distinction matters: graph couplings can approximate general couplings when
the source is nonatomic, but the approximating maps need not converge to a
map.

\subsection{Power-distance GW and its Monge restriction}

Let \(\mathcal X\subset\R^n\) and \(\mathcal Y\subset\R^d\) be compact, let
\(\alpha\in\Pcal(\mathcal X)\) and \(\beta\in\Pcal(\mathcal Y)\), and fix
\(0<p\leq2\).  Write
\[
  d_{p,\mathcal X}(x,x')=\norm{x-x'}^p,
  \qquad
  d_{p,\mathcal Y}(y,y')=\norm{y-y'}^p.
\]
The inner exponent \(p\) should not be confused with the outer quadratic
distortion, which remains fixed throughout this note.  Thus \(p=1\) is the
usual quadratic GW problem built from Euclidean distances, whereas \(p=2\)
is the squared-Euclidean-dissimilarity model studied in
\citet{DumontLacombeVialard2025}.

\begin{definition}[Power-distance GW problem]
\label{def:power-gw}
For \(\pi\in\Piopt(\alpha,\beta)\), set
\begin{equation}
  \Ecal_p(\pi)
  :=
  \iint
  \bigl(d_{p,\mathcal X}(x,x')-d_{p,\mathcal Y}(y,y')\bigr)^2
  \dd\pi(x,y)\dd\pi(x',y').
  \label{eq:power-gw-energy}
\end{equation}
The Kantorovich and Monge values are
\begin{align}
  \GW_{p,2}^2(\alpha,\beta)
  &:=\min_{\pi\in\Piopt(\alpha,\beta)}\Ecal_p(\pi),
  \label{eq:power-gw-kantorovich}\\
  \GM_{p,2}^2(\alpha,\beta)
  &:=\inf_{T_\sharp\alpha=\beta}
  \Ecal_p\bigl((\Id,T)_\sharp\alpha\bigr),
  \label{eq:power-gw-monge}
\end{align}
with value \(+\infty\) in \eqref{eq:power-gw-monge} if no admissible map
exists.
\end{definition}

Compactness makes existence on the Kantorovich side immediate.

\begin{proposition}[Existence of a coupling minimizer]
\label{prop:gw-existence}
The minimum in \eqref{eq:power-gw-kantorovich} is attained.
\end{proposition}

\begin{proof}
The set \(\Piopt(\alpha,\beta)\) is narrowly compact.  The integrand in
\eqref{eq:power-gw-energy} is continuous and bounded on
\((\mathcal X\times\mathcal Y)^2\).  Hence
\(\pi\mapsto\Ecal_p(\pi)\) is narrowly continuous.
\end{proof}

The central question is whether some minimizer in
\eqref{eq:power-gw-kantorovich} has the form \((\Id,T)_\sharp\alpha\).  This
is stronger than equality of the two infimum values.  The distinction follows
from the classical density of deterministic couplings
\citep{Pratelli2007,MemoliNeedham2024}, which takes the following particularly
transparent form in the present compact setting.

\begin{proposition}[Graph couplings are dense, but need not be closed]
\label{prop:graph-couplings-dense}
Assume that \(\alpha\) is nonatomic.  For every
\(\pi\in\Piopt(\alpha,\beta)\), there are measurable maps
\(T_k:\mathcal X\to\mathcal Y\) such that
\begin{equation}
  (T_k)_\sharp\alpha=\beta,
  \qquad
  (\Id,T_k)_\sharp\alpha\rightharpoonup\pi.
  \label{eq:dense-graph-couplings}
\end{equation}
Consequently, for every \(0<p\leq2\),
\begin{equation}
  \GM_{p,2}^2(\alpha,\beta)=\GW_{p,2}^2(\alpha,\beta).
  \label{eq:gm-gw-value-equality}
\end{equation}
Neither the convergence in \eqref{eq:dense-graph-couplings} nor the equality
in \eqref{eq:gm-gw-value-equality} implies that the Gromov--Monge infimum is
attained.
\end{proposition}

\begin{proof}
Choose finite Borel partitions
\(\mathcal X=\bigsqcup_i A_i^k\) whose cells have diameter at most \(1/k\),
and define finite measures on \(\mathcal Y\) by
\begin{equation*}
  \beta_i^k(B):=\pi(A_i^k\times B).
\end{equation*}
The measures \(\alpha|_{A_i^k}\) and \(\beta_i^k\) have the same mass.  Since
the former is nonatomic, the standard mapping theorem for probability
measures on standard Borel spaces
\citep[Theorem~2.1]{DumontLacombeVialard2025} gives, on every cell of positive
mass, a measurable map \(T_i^k:A_i^k\to\mathcal Y\) satisfying
\((T_i^k)_\sharp(\alpha|_{A_i^k})=\beta_i^k\).  Joining these finitely many
maps defines \(T_k\), and summing over the cells gives
\((T_k)_\sharp\alpha=\sum_i\beta_i^k=\beta\).

Let \(q\in C(\mathcal X\times\mathcal Y)\) and let \(\omega_q\) be its
uniform modulus of continuity in the first variable.  Choosing
\(x_i^k\in A_i^k\), both the restriction of
\((\Id,T_k)_\sharp\alpha\) to \(A_i^k\times\mathcal Y\) and that of \(\pi\)
have target marginal \(\beta_i^k\).  Replacing their first coordinate by
\(x_i^k\) therefore shows that
\begin{equation*}
  \left|\int q\dd(\Id,T_k)_\sharp\alpha-\int q\dd\pi\right|
  \leq 2\omega_q(1/k)\longrightarrow0.
\end{equation*}
This proves \eqref{eq:dense-graph-couplings}.  Products of narrowly convergent
probability measures converge narrowly on the compact product space, so the
continuity of the integrand in \eqref{eq:power-gw-energy} yields
\(\Ecal_p((\Id,T_k)_\sharp\alpha)\to\Ecal_p(\pi)\).  Applying this to a GW
minimizer proves the inequality \(\GM_{p,2}^2\leq\GW_{p,2}^2\); the reverse
inequality follows because every graph coupling is admissible for GW.
Finally, if \(\beta\) is not a Dirac mass, the product coupling
\(\alpha\otimes\beta\) is not graph-supported, although the first part of the
proposition approximates it by graph couplings.  Thus graph couplings are not
narrowly closed in general.
\end{proof}

Only the source is required to be nonatomic in
\cref{prop:graph-couplings-dense}; the target may be discrete.  The symmetric
metric-space comparison of \citet{MemoliNeedham2024} treats substantially more
general, possibly unbounded network functions.  Here compactness and
continuity make the approximation argument elementary.  The Brenier problem
studied below is precisely the additional question of whether the dense class
of graph couplings contains an actual GW minimizer.

\subsection{What is known at \texorpdfstring{\(p=2\)}{p=2}}

Classical Brenier theory concerns a \emph{linear} transport objective with
quadratic ground cost: if the source has a density, the unique optimal
coupling is induced by the gradient of a convex function
\citep{Brenier1991,Villani2009,Santambrogio2015}.  GW is instead quadratic in
the unknown coupling, so Brenier's theorem cannot be applied before one has
identified an appropriate frozen linear cost.

For GW built from inner products, this program succeeds particularly well.
Vayer proved a transformed-Brenier theorem under a full-rank assumption
\citep[Theorem~4.2.3]{Vayer2020}, and
\citet{DumontLacombeVialard2025} subsequently established existence of a
Monge optimizer without that rank assumption.  Squared Euclidean
dissimilarities behave differently.  In the notation of
\cref{def:power-gw}, this is exactly the case \(p=2\).

A foundational positive result is available under symmetry.  Sturm proved
that when two absolutely continuous measures on the same \(\R^d\) are
rotationally invariant around their barycenters, every squared-distance GW
optimizer is a radial quantile rearrangement followed by an orthogonal map
\citep[Theorem~9.21]{Sturm2023Space}.  We restate the result and unpack its
polar proof in \cref{thm:sturm-radial-brenier}; importantly, its assumptions
involve only the marginals, not an unknown optimal coupling.

\begin{definition}[Finite-map coupling]
\label{def:two-map}
A coupling \(\pi\in\Piopt(\alpha,\beta)\) is a \emph{\(k\)-map} if there are
measurable maps \(T_1,\ldots,T_k:\mathcal X\to\mathcal Y\) and measurable
weights \(\lambda_j:\mathcal X\to[0,1]\), with
\(\sum_{j=1}^k\lambda_j=1\) \(\alpha\)-almost everywhere, such that
\[
  \pi=\sum_{j=1}^k(\Id,T_j)_\sharp(\lambda_j\alpha).
\]
Equivalently, the conditional measure \(\pi_x\) is supported on at most \(k\)
points for \(\alpha\)-almost every \(x\).  A \emph{two-map} is the case
\(k=2\).
\end{definition}

A different two-limb structure is naturally selected by the subtwist
condition.  Unlike a two-map, its second limb is parameterized from the target
side.

\begin{definition}[Map/anti-map coupling]
\label{def:map-antimap}
A coupling \(\pi\in\Piopt(\alpha,\beta)\) is a \emph{map/anti-map} if there
are a submeasure \(\bar\alpha\leq\alpha\) and measurable maps
\(G:\mathcal X\to\mathcal Y\), \(H:\mathcal Y\to\mathcal X\) such that
\begin{equation}
  \pi=(\Id,G)_\sharp\bar\alpha
      +(H,\Id)_\sharp\bigl(\beta-G_\sharp\bar\alpha\bigr).
  \label{eq:map-antimap-decomposition}
\end{equation}
Here \(G_\sharp\bar\alpha\leq\beta\), the second target submeasure gives no
mass to the range of \(G\), and its image by \(H\) is
\(\alpha-\bar\alpha\).  Thus the support is contained in one graph from
source to target and one anti-graph from target to source.
\end{definition}

The following is the current general structural theorem for squared-distance
GW.

\begin{theorem}[Known squared-distance structure
  \citealp{DumontLacombeVialard2025}]
\label{thm:known-two-map}
Let \(n\geq d\), let \(\alpha\) and \(\beta\) have compact support, and assume
\(\alpha\ll\mathrm{Leb}^n\).  For \(p=2\), an optimal two-map exists.  More
precisely, if
\[
  S_\pi:=\int
  (y-m_\beta)(x-m_\alpha)^\top\dd\pi(x,y)
  \in\R^{d\times n}
\]
is the cross-covariance of an optimal coupling, their theorem gives a
single-map optimizer when \(\rank S_\pi\leq d-2\), and a two-map in the
remaining rank regimes.  In the full-rank regime the frozen linear optimizer
is unique and has both the two-map and map/anti-map structures of
\cref{def:two-map,def:map-antimap}.
\end{theorem}

The same work gives one-dimensional numerical examples in which the optimum
appears genuinely two-valued.  This is evidence, rather than a theorem of
nonexistence, but it makes a universal single-map claim mathematically
unsafe.  Our objective below is therefore twofold: derive the strongest
general theorem furnished by the RKHS reduction, and identify concrete
hypotheses under which either a frozen optimizer has finite twist or a
twisted quotient problem can be lifted to a Monge optimizer.

There is nevertheless a sharp positive conclusion once the source has more
ambient dimensions than the target.  The rank alternatives in the preceding
theorem can then be closed: low-rank costs have atomless source fibers that can
be used for a deterministic lift, while full-row-rank costs are twisted away
from a lower-dimensional exceptional set.  The resulting rank-free theorem
is proved in \cref{thm:p2-strict-dimension-gap-monge}.

\section{Power distances as Hilbert features}
\label{sec:hilbert-features}

The restriction \(0<p\leq2\) is exactly the range in which Euclidean power
distances are of negative type.  This converts the nonlinear-looking
distance product in GW into a squared Hilbert--Schmidt norm.

\subsection{Schoenberg embedding and centered kernels}

\begin{proposition}[Schoenberg feature representation]
\label{prop:schoenberg-feature}
For every \(0<p\leq2\) and every Euclidean dimension \(m\), there are a real
Hilbert space \(\Hcal_{p,m}\) and an injective continuous map
\(\Phi_{p,m}:\R^m\to\Hcal_{p,m}\) such that
\begin{equation}
  \norm{x-x'}^p
  =\norm{\Phi_{p,m}(x)-\Phi_{p,m}(x')}^2_{\Hcal_{p,m}}.
  \label{eq:schoenberg-embedding}
\end{equation}
For \(p=2\), one can take \(\Hcal_{2,m}=\R^m\) and
\(\Phi_{2,m}=\Id\).  For \(0<p<2\), one possible realization is the real
form of
\begin{equation}
  \Phi_{p,m}(x)(\omega)
  =c_{m,p}^{1/2}\bigl(e^{\mathrm i\ip{\omega}{x}}-1\bigr)
  \quad\text{in}\quad
  L^2\!\left(\R^m,\frac{\dd\omega}{\norm{\omega}^{m+p}}\right),
  \label{eq:fourier-schoenberg-feature}
\end{equation}
for a normalization constant \(c_{m,p}>0\).
\end{proposition}

\begin{proof}
Schoenberg's theorem says that \((x,x')\mapsto\norm{x-x'}^p\) is
conditionally negative definite precisely in the range \(0<p\leq2\); the
standard Hilbert embedding of a negative-type semimetric gives
\eqref{eq:schoenberg-embedding} \citep{Schoenberg1938}.  For \(0<p<2\), the
identity
\[
  \norm{x-x'}^p
  =c_{m,p}\int_{\R^m}
  \frac{\abs{e^{\mathrm i\ip{\omega}{x}}
             -e^{\mathrm i\ip{\omega}{x'}}}^2}
       {\norm{\omega}^{m+p}}\dd\omega
\]
gives \eqref{eq:fourier-schoenberg-feature}.  The integral is finite near
the origin because \(p<2\), and at infinity because \(p>0\).  The identity
map proves the endpoint \(p=2\).
\end{proof}

Apply \cref{prop:schoenberg-feature} separately to \(\mathcal X\) and
\(\mathcal Y\), writing the feature maps as \(\Phi\) and \(\Psi\).  Center
them with respect to the marginals:
\begin{equation}
  u(x):=\Phi(x)-\int\Phi\dd\alpha,
  \qquad
  v(y):=\Psi(y)-\int\Psi\dd\beta.
  \label{eq:centered-features}
\end{equation}
Let \(\Hcal_\alpha\) and \(\Hcal_\beta\) be the closed linear spans of
\(u(\mathcal X)\) and \(v(\mathcal Y)\).  Their reproducing kernels are
\begin{align}
  k_\alpha(x,x')&:=\ip{u(x)}{u(x')}_{\Hcal_\alpha},
  &
  k_\beta(y,y')&:=\ip{v(y)}{v(y')}_{\Hcal_\beta}.
  \label{eq:centered-power-kernels}
\end{align}
These kernels can be evaluated without constructing the feature spaces.  Set
\begin{align}
  r_\alpha(x)&:=\int d_{p,\mathcal X}(x,z)\dd\alpha(z),
  &
  D_\alpha&:=\iint d_{p,\mathcal X}(x,z)\dd\alpha(x)\dd\alpha(z),
  \label{eq:distance-means-alpha}\\
  r_\beta(y)&:=\int d_{p,\mathcal Y}(y,z)\dd\beta(z),
  &
  D_\beta&:=\iint d_{p,\mathcal Y}(y,z)\dd\beta(y)\dd\beta(z).
  \label{eq:distance-means-beta}
\end{align}
Double centering gives
\begin{align}
  k_\alpha(x,x')
  &=\frac12\bigl(r_\alpha(x)+r_\alpha(x')
    -d_{p,\mathcal X}(x,x')-D_\alpha\bigr),
  \label{eq:centered-kernel-alpha-explicit}\\
  k_\beta(y,y')
  &=\frac12\bigl(r_\beta(y)+r_\beta(y')
    -d_{p,\mathcal Y}(y,y')-D_\beta\bigr).
  \label{eq:centered-kernel-beta-explicit}
\end{align}
In particular, define
\begin{equation}
  a(x):=k_\alpha(x,x)=\norm{u(x)}^2,
  \qquad
  b(y):=k_\beta(y,y)=\norm{v(y)}^2.
  \label{eq:feature-radii}
\end{equation}
Their means are \(\bar a=D_\alpha/2\) and \(\bar b=D_\beta/2\).

\subsection{The cross-covariance decomposition}

For \(\pi\in\Piopt(\alpha,\beta)\), define the Hilbert--Schmidt operator
\begin{equation}
  C_\pi
  :=\int v(y)\otimes u(x)\dd\pi(x,y)
  \in\HS(\Hcal_\alpha,\Hcal_\beta),
  \label{eq:hilbert-cross-covariance}
\end{equation}
where \((v\otimes u)h=\ip{u}{h}v\).  Compactness makes this Bochner integral
well defined.

\begin{proposition}[Distance-product decomposition]
\label{prop:distance-product-decomposition}
For every \(\pi\in\Piopt(\alpha,\beta)\),
\begin{equation}
  \iint d_{p,\mathcal X}(x,x')d_{p,\mathcal Y}(y,y')
  \dd\pi(x,y)\dd\pi(x',y')
  =2\int a(x)b(y)\dd\pi(x,y)
   +2\bar a\bar b+4\norm{C_\pi}_{\HS}^2.
  \label{eq:distance-product-decomposition}
\end{equation}
Consequently,
\begin{equation}
  \Ecal_p(\pi)
  =\mathsf C_p(\alpha,\beta)
   -4\int a(x)b(y)\dd\pi(x,y)
   -8\norm{C_\pi}_{\HS}^2,
  \label{eq:gw-feature-decomposition}
\end{equation}
where the marginal-only constant is
\begin{equation}
  \mathsf C_p(\alpha,\beta)
  :=\iint d_{p,\mathcal X}(x,x')^2\dd\alpha(x)\dd\alpha(x')
    +\iint d_{p,\mathcal Y}(y,y')^2\dd\beta(y)\dd\beta(y')
    -4\bar a\bar b.
  \label{eq:gw-marginal-constant}
\end{equation}
In particular, \(\pi\mapsto\Ecal_p(\pi)\) is concave on
\(\Piopt(\alpha,\beta)\).
\end{proposition}

\begin{proof}
Let \((X,Y)\) and \((X',Y')\) be independent with law \(\pi\).  By
\eqref{eq:schoenberg-embedding} and \eqref{eq:centered-features},
\[
  d_{p,\mathcal X}(X,X')
  =a(X)+a(X')-2\ip{u(X)}{u(X')},
\]
and similarly on \(\mathcal Y\).  Expand the product.  The product of the
two radial sums contributes
\(2\mathbb E[a(X)b(Y)]+2\bar a\bar b\).  Every mixed term containing one
inner product vanishes because \(\mathbb E[u(X)]=0\) and
\(\mathbb E[v(Y)]=0\).  Finally,
\begin{align*}
  \mathbb E\!\left[
    \ip{u(X)}{u(X')}\ip{v(Y)}{v(Y')}
  \right]
  &=\ip{C_\pi}{C_\pi}_{\HS}
   =\norm{C_\pi}_{\HS}^2.
\end{align*}
This proves \eqref{eq:distance-product-decomposition}.  Expanding the square
in \eqref{eq:power-gw-energy} then gives
\eqref{eq:gw-feature-decomposition}.  Its last term is the negative of the
square of a linear function of \(\pi\), while the preceding coupling term is
linear; hence the objective is concave.
\end{proof}

\begin{corollary}[Discrete permutation optimizer]
\label{cor:finite-permutation-optimizer}
Suppose
\[
 \alpha=N^{-1}\sum_{i=1}^N\delta_{x_i},
 \qquad
 \beta=N^{-1}\sum_{j=1}^N\delta_{y_j}.
\]
For every \(0<p\leq2\), the problem \eqref{eq:power-gw-kantorovich} has an
optimal permutation coupling
\begin{equation}
  \pi_\sigma=\frac1N\sum_{i=1}^N\delta_{(x_i,y_{\sigma(i)})}
  \label{eq:finite-permutation-coupling}
\end{equation}
for some permutation \(\sigma\) of \(\{1,\ldots,N\}\).
\end{corollary}

\begin{proof}
The coupling set is the scaled Birkhoff polytope.  Represent any minimizer as
a convex combination of its extreme points.  Concavity and minimality imply
that every extreme point carrying positive weight is also a minimizer.  The
Birkhoff--von Neumann theorem identifies these extreme points with the
permutation couplings \eqref{eq:finite-permutation-coupling}.
\end{proof}

This is a genuine discrete Monge statement, but its proof does not pass to
continuous measures: extreme couplings over nonatomic spaces need not be
graphs, and the class of graph couplings is not narrowly closed.  Concavity
is therefore useful structurally without making the continuous minimization
easy.  See \citet{Vayer2026Sparsity} for a recent account of this extremal
mechanism and related sparsity properties of GW plans.  The feature
decomposition becomes more informative after dualizing the squared
Hilbert--Schmidt norm.

\section{An exact RKHS-regularized learned-cost problem}
\label{sec:learned-cost}

The tensor-product space
\(
 \Hcal_\beta\widehat\otimes\Hcal_\alpha
\)
is canonically identified with
\(\HS(\Hcal_\alpha,\Hcal_\beta)\).  Under this identification, an operator
\(A\) defines the cross-space function
\begin{equation}
  h_A(x,y):=\ip{Au(x)}{v(y)}_{\Hcal_\beta}.
  \label{eq:tensor-rkhs-function}
\end{equation}
It belongs to the RKHS with product kernel
\begin{equation}
  K_{\alpha,\beta}\bigl((x,y),(x',y')\bigr)
  =k_\alpha(x,x')k_\beta(y,y'),
  \qquad
  \norm{h_A}_{K_{\alpha,\beta}}=\norm{A}_{\HS}.
  \label{eq:tensor-rkhs-kernel}
\end{equation}
The equality of norms uses the definition of \(\Hcal_\alpha\) and
\(\Hcal_\beta\) as the closed feature spans.

For a continuous cost \(c\) on \(\mathcal X\times\mathcal Y\), write
\begin{equation}
  \Lcal_c(\alpha,\beta)
  :=\min_{\pi\in\Piopt(\alpha,\beta)}\int c(x,y)\dd\pi(x,y).
  \label{eq:linear-ot-value}
\end{equation}

\begin{theorem}[Exact RKHS cost learning]
\label{thm:exact-rkhs-cost-learning}
For \(A\in\HS(\Hcal_\alpha,\Hcal_\beta)\), define
\begin{equation}
  c_A(x,y):=-4a(x)b(y)-4\ip{Au(x)}{v(y)}.
  \label{eq:learned-cost}
\end{equation}
Then, for every coupling \(\pi\),
\begin{equation}
  \Ecal_p(\pi)
  =\mathsf C_p(\alpha,\beta)
   +\min_A\left\{
      \frac12\norm{A}_{\HS}^2+
      \int c_A\dd\pi
    \right\},
  \qquad
  A_\pi=4C_\pi.
  \label{eq:pointwise-cost-learning}
\end{equation}
Consequently,
\begin{equation}
  \GW_{p,2}^2(\alpha,\beta)
  =\mathsf C_p(\alpha,\beta)
   +\min_{A\in\HS(\Hcal_\alpha,\Hcal_\beta)}
    \left\{\frac12\norm{A}_{\HS}^2+
    \Lcal_{c_A}(\alpha,\beta)\right\}.
  \label{eq:global-cost-learning}
\end{equation}
Equivalently, with
\(
 r_A(x,y)=4a(x)b(y)+4\ip{Au(x)}{v(y)},
\)
\begin{equation}
  \GW_{p,2}^2(\alpha,\beta)
  =\mathsf C_p(\alpha,\beta)
  -\max_{\substack{A\in\HS(\Hcal_\alpha,\Hcal_\beta)\\
                     \pi\in\Piopt(\alpha,\beta)}}
   \left\{
     \int r_A\dd\pi-\frac12\norm{A}_{\HS}^2
   \right\}.
  \label{eq:regularized-max-max}
\end{equation}
The maximand is strictly concave in \(A\) for fixed \(\pi\) and linear in
\(\pi\) for fixed \(A\); it is not asserted to be jointly concave.
\end{theorem}

\begin{proof}
The Hilbert--Schmidt pairing satisfies
\[
  \ip{A}{C_\pi}_{\HS}
  =\int\ip{Au(x)}{v(y)}\dd\pi(x,y).
\]
Completing the square gives
\[
  \min_A\left\{
    \frac12\norm{A}_{\HS}^2-4\ip{A}{C_\pi}_{\HS}
  \right\}
  =-8\norm{C_\pi}_{\HS}^2,
  \qquad A_\pi=4C_\pi.
\]
Together with \eqref{eq:gw-feature-decomposition}, this proves
\eqref{eq:pointwise-cost-learning}.  Taking the joint minimum in \((\pi,A)\)
and minimizing first in \(\pi\) proves \eqref{eq:global-cost-learning}.
Changing signs and using \(c_A=-r_A\) gives
\eqref{eq:regularized-max-max}.
\end{proof}

This is the promised RKHS regularizer.  The fixed radial product
\(-4a(x)b(y)\) is not learned; the learned component lies in the
tensor-product RKHS and is penalized by its squared norm.  The representation
is closely related to cost-regularized cross-space transport
\citep{SebbouhCuturiPeyre2024} and to negative-type biconvex GW relaxations
\citep{Konno1976,SejourneVialardPeyre2021}.

\begin{remark}[A bilinear cost on augmented features]
\label{rem:augmented-feature-cost}
Introduce
\[
  \widehat u(x):=(a(x),u(x))\in\R\oplus\Hcal_\alpha,
  \qquad
  \widehat v(y):=(b(y),v(y))\in\R\oplus\Hcal_\beta,
\]
and the block operator
\(
 \widehat A(s,h):=(s,Ah).
\)
Then the learned cost is the bilinear expression
\begin{equation}
  c_A(x,y)=-4\ip{\widehat A\widehat u(x)}{\widehat v(y)}.
  \label{eq:augmented-bilinear-cost}
\end{equation}
Only the block \(A\) is optimized and regularized; the scalar block is fixed.
Thus the reduction is bilinear after a nonlinear feature lift, not in the
original variables.  For \(p=2\), these augmented features are
\(
 (\norm{x-m_\alpha}^2,x-m_\alpha)
\)
and
\(
 (\norm{y-m_\beta}^2,y-m_\beta)
\),
which lie on paraboloids.  The radial coordinate of the target paraboloid is
exactly what later creates a possible second branch.
\end{remark}

\begin{corollary}[Self-consistent frozen OT problem]
\label{cor:self-consistent-linear-ot}
A pair \((\pi_\star,A_\star)\) minimizes the joint problem in
\eqref{eq:global-cost-learning} if and only if \(\pi_\star\) is GW-optimal,
\begin{equation}
  A_\star=4C_{\pi_\star},
  \label{eq:self-consistent-operator}
\end{equation}
and \(\pi_\star\) solves ordinary Kantorovich transport for the cost
\(c_{A_\star}\).  At self-consistency this cost can be evaluated only from
the centered kernels:
\begin{equation}
  c_{\pi_\star}(x,y)
  =-4a(x)b(y)
   -16\int k_\alpha(x,x')k_\beta(y,y')
     \dd\pi_\star(x',y').
  \label{eq:self-consistent-kernel-cost}
\end{equation}
\end{corollary}

\begin{proof}
If \((\pi_\star,A_\star)\) is jointly optimal, minimizing first in \(A\) and
using \eqref{eq:pointwise-cost-learning} gives
\(A_\star=4C_{\pi_\star}\) and shows that \(\pi_\star\) is GW-optimal.  With
\(A_\star\) frozen, replacing \(\pi_\star\) by a better linear-OT coupling
would decrease the joint objective, so \(\pi_\star\) solves that frozen
problem.  Conversely, if \(\pi_\star\) is GW-optimal and
\(A_\star=4C_{\pi_\star}\), then \eqref{eq:pointwise-cost-learning} makes the
pair jointly optimal; its frozen optimality follows by the preceding
replacement argument.  Finally,
\[
  \ip{C_{\pi_\star}u(x)}{v(y)}
  =\int k_\alpha(x,x')k_\beta(y,y')\dd\pi_\star(x',y'),
\]
which gives \eqref{eq:self-consistent-kernel-cost}.
\end{proof}

The corollary is the bridge to Monge theory: the nonlinear GW minimizer is a
linear OT minimizer once its learned cost has been frozen.  The remaining
question is whether this particular cost forces graph support.

The freezing principle is considerably more general than the negative-type
feature representation.  Let \(Q\) be any continuous symmetric kernel on
\((\mathcal X\times\mathcal Y)^2\), and define
\begin{equation}
  \Ecal_Q(\pi):=\iint Q(z,z')\dd\pi(z)\dd\pi(z'),
  \qquad z=(x,y).
  \label{eq:abstract-quadratic-gw}
\end{equation}
This includes a general GW distortion by taking
\(
 Q((x,y),(x',y'))=L(d_{\mathcal X}(x,x'),d_{\mathcal Y}(y,y')).
\)

\begin{proposition}[Universal frozen-cost and replacement principles]
\label{prop:universal-frozen-gw}
Let \(\pi_\star\) minimize \(\Ecal_Q\) over \(\Piopt(\alpha,\beta)\), and set
\begin{equation}
  c_{\pi_\star}(z)
  :=2\int Q(z,z')\dd\pi_\star(z').
  \label{eq:abstract-frozen-gw-cost}
\end{equation}
Then \(\pi_\star\) minimizes the linear OT problem with cost
\(c_{\pi_\star}\).  This conclusion requires no convexity or concavity of
\(\Ecal_Q\).

More precisely, if \(\widetilde\pi\) is any minimizer of the frozen linear
problem and \(\eta=\widetilde\pi-\pi_\star\), then
\begin{equation}
  \Ecal_Q(\widetilde\pi)-\Ecal_Q(\pi_\star)
  =\iint Q(z,z')\dd\eta(z)\dd\eta(z')\geq0.
  \label{eq:frozen-replacement-defect}
\end{equation}
Consequently, \(\widetilde\pi\) is GW-optimal if and only if the quadratic
term in \eqref{eq:frozen-replacement-defect} vanishes.

Suppose, in addition, that
\begin{equation}
  \iint Q(z,z')\dd\eta(z)\dd\eta(z')\leq0
  \quad\text{for }\eta=\pi_1-\pi_0,
  \quad \pi_0,\pi_1\in\Piopt(\alpha,\beta).
  \label{eq:conditional-concavity-gw-kernel}
\end{equation}
Equivalently, \(\Ecal_Q\) is concave on
\(\Piopt(\alpha,\beta)\).  Then every minimizer of the frozen linear problem
for \(c_{\pi_\star}\) is also a global minimizer of \(\Ecal_Q\).
It is sufficient, but generally stronger, to require the same inequality for
every finite signed \(\eta\) with zero \(\mathcal X\)- and
\(\mathcal Y\)-marginals.
\end{proposition}

\begin{proof}
For \(\pi\in\Piopt(\alpha,\beta)\), put \(\eta=\pi-\pi_\star\).  Expanding
along the feasible segment gives
\begin{equation}
  \Ecal_Q(\pi_\star+t\eta)
  =\Ecal_Q(\pi_\star)
   +t\int c_{\pi_\star}\dd\eta
   +t^2\iint Q\dd\eta\dd\eta.
  \label{eq:quadratic-gw-segment-expansion}
\end{equation}
Minimality at \(t=0\) forces
\(\int c_{\pi_\star}\dd\eta\geq0\), proving frozen optimality.  If
\(\widetilde\pi\) is another frozen minimizer, this linear term vanishes for
\(\eta=\widetilde\pi-\pi_\star\).  Setting \(t=1\) proves the equality in
\eqref{eq:frozen-replacement-defect}; its inequality follows from global
minimality of \(\pi_\star\).  Under
\eqref{eq:conditional-concavity-gw-kernel}, the expansion at \(t=1\) gives
\(\Ecal_Q(\widetilde\pi)\leq\Ecal_Q(\pi_\star)\).  The reverse inequality
follows from global minimality of \(\pi_\star\), hence equality holds.
\end{proof}

The first assertion is a necessary first-order condition for every quadratic
GW problem.  The second is stronger: conditional concavity permits replacing
an optimizer by a more structured optimizer of its frozen problem.  This is
the mechanism needed below to manufacture a Monge GW optimizer through a
quotient, rather than merely diagnose whether a prescribed optimizer is a
map.  For power-distance GW, conditional concavity is exactly
\eqref{eq:gw-feature-decomposition}.

The same frozen problem admits a direct expression that does not require
constructing the feature spaces.  It is the first variation of the original
quadratic GW objective after discarding terms whose integral is fixed by the
marginals.

\begin{proposition}[Direct first-variation cost]
\label{prop:direct-first-variation-cost}
For \(\pi\in\Piopt(\alpha,\beta)\), define
\begin{equation}
  \ell_\pi(x,y)
  :=-4\int
  d_{p,\mathcal X}(x,x')d_{p,\mathcal Y}(y,y')
  \dd\pi(x',y').
  \label{eq:direct-first-variation-cost}
\end{equation}
If \(A_\pi=4C_\pi\), then there are continuous functions
\(\varphi_\pi:\mathcal X\to\R\), \(\psi_\pi:\mathcal Y\to\R\), and a
constant \(q_\pi\) such that
\begin{equation}
  \ell_\pi(x,y)
  =c_{A_\pi}(x,y)+\varphi_\pi(x)+\psi_\pi(y)+q_\pi.
  \label{eq:equivalent-frozen-costs}
\end{equation}
Consequently, \(\ell_\pi\) and \(c_{A_\pi}\) have the same optimal
couplings.  Wherever both costs are differentiable in \(x\), their twist maps
differ by a vector independent of \(y\), and hence have the same injectivity.
In particular, every GW minimizer \(\pi_\star\) solves linear OT for
\(\ell_{\pi_\star}\).

If \(1<p\leq2\), differentiation under the integral gives
\begin{equation}
  \nabla_x\ell_\pi(x,y)
  =-4p\int
  \norm{x-x'}^{p-2}(x-x')\norm{y-y'}^p
  \dd\pi(x',y'),
  \label{eq:gradient-direct-first-variation-cost}
\end{equation}
where the vector integrand is set to zero at \(x=x'\).
\end{proposition}

\begin{proof}
For a feasible signed direction \(\eta\) with zero marginals, the first
variation is
\[
  D\Ecal_p(\pi)[\eta]
  =2\iint
  \bigl(d_{p,\mathcal X}(x,x')-d_{p,\mathcal Y}(y,y')\bigr)^2
  \dd\pi(x',y')\dd\eta(x,y).
\]
Thus it is represented on the fixed-marginal affine set by the function
\[
  (x,y)\longmapsto
  2\int
  \bigl(d_{p,\mathcal X}(x,x')-d_{p,\mathcal Y}(y,y')\bigr)^2
  \dd\pi(x',y').
\]
The two squared terms depend only on \(x\) or only on \(y\); their integrals
are fixed on \(\Piopt(\alpha,\beta)\).  Removing them leaves
\eqref{eq:direct-first-variation-cost}.  Alternatively, insert
\[
  d_{p,\mathcal X}(x,x')=a(x)+a(x')-2k_\alpha(x,x')
\]
and its \(\mathcal Y\)-counterpart into
\eqref{eq:direct-first-variation-cost}.  The only terms depending jointly on
\((x,y)\) are
\[
  -4a(x)b(y)
  -16\int k_\alpha(x,x')k_\beta(y,y')\dd\pi(x',y'),
\]
which equal \(c_{A_\pi}\) by
\eqref{eq:self-consistent-kernel-cost}; all remaining terms are separable or
constant.  This proves \eqref{eq:equivalent-frozen-costs}.  Such additions do
not change a Kantorovich optimizer, and their \(x\)-gradients do not depend on
\(y\), so they do not change twist.  The last assertion follows either from
\cref{cor:self-consistent-linear-ot} or by applying
\cref{prop:universal-frozen-gw} to the quadratic functional \(\Ecal_p\).
Formula
\eqref{eq:gradient-direct-first-variation-cost} follows by differentiating
\(\norm{x-x'}^p\).
\end{proof}

\section{The Brenier mechanism: twist of the learned cost}
\label{sec:rkhs-twist}

For general costs, the appropriate replacement for strict convexity of the
quadratic cost is control of the twist fibers.  One branch gives a Monge map;
finitely many branches give a finite mixture of maps.  Since the source is
absolutely continuous, these conditions are needed only almost everywhere in
the source variable and only on the dual contact set.

\subsection{Finite twist and finite-map solutions}

For a dual feasible pair \((f,g)\), define its contact set
\begin{equation}
  S_{f,g}:=\{(x,y):f(x)+g(y)=c(x,y)\}.
  \label{eq:dual-contact-set}
\end{equation}
Every optimal plan is concentrated on this set when \((f,g)\) is dual
optimal.  It is therefore unnecessary to control all fibers of the twist
map; controlling the fibers that meet \(S_{f,g}\) is enough.

\begin{definition}[Finite twist on a contact set]
\label{def:ae-twist}
Let \(c\) be differentiable in its first variable and let
\(S\subset\mathcal X\times\mathcal Y\).  The cost is
\emph{\(m\)-twisted on \(S\) at \(x\)} if, for every \(q\in\R^n\),
\begin{equation}
  \card\{y\in\supp\beta:(x,y)\in S,
            \ \nabla_xc(x,y)=q\}\leq m.
  \label{eq:m-twist-contact}
\end{equation}
It is \emph{\(\alpha\)-almost-everywhere \(m\)-twisted on \(S\)} if
\eqref{eq:m-twist-contact} holds for \(\alpha\)-almost every \(x\).  The
global \(m\)-twist condition is obtained by taking
\(S=\supp\alpha\times\supp\beta\); ordinary twist is the case \(m=1\).
\end{definition}

This finite-fiber relaxation of twist is precisely the condition that yields
a finite mixture of transport maps \citep{MoameniPass2017}.

\begin{proposition}[Finite twist selects finitely many graphs]
\label{prop:finite-twist-linear-ot}
Assume \(\alpha\ll\mathrm{Leb}^n\), the supports are compact, and \(c\) is
continuous and \(C^1\) in its first variable near the supports.  Let
\((f,g)\) be a dual optimizer for \eqref{eq:linear-ot-value}.  If \(c\) is
\(\alpha\)-almost-everywhere \(m\)-twisted on \(S_{f,g}\), then every primal
optimizer \(\pi\) is a \(k\)-map for some \(k\leq m\): there are measurable
maps \(T_j\) and weights \(\lambda_j\geq0\),
\(\sum_{j=1}^k\lambda_j=1\), such that
\begin{equation}
  \pi=\sum_{j=1}^k(\Id,T_j)_\sharp(\lambda_j\alpha).
  \label{eq:finite-map-decomposition}
\end{equation}
At every differentiability point of \(f\), each active branch satisfies
\begin{equation}
  \nabla f(x)=\nabla_xc(x,T_j(x))
  \qquad\text{whenever }\lambda_j(x)>0.
  \label{eq:finite-branch-characterization}
\end{equation}
\end{proposition}

\begin{proof}
The uniform bound on \(\nabla_xc\) makes the \(c\)-concave potential \(f\)
Lipschitz, hence differentiable \(\alpha\)-almost everywhere.  Disintegrate
an optimal plan as \(\pi=\int\delta_x\otimes\pi_x\dd\alpha(x)\).  At a
differentiability point \(x\), every \(y\in\supp\pi_x\) lies in the contact
set.  Since \(f(z)\leq c(z,y)-g(y)\), with equality at \(z=x\), differentiation
gives
\(
 \nabla f(x)=\nabla_xc(x,y).
\)
Condition \eqref{eq:m-twist-contact} therefore implies
\(\card(\supp\pi_x)\leq m\).  A finite-valued Borel multifunction admits a
measurable enumeration; enumerating the atoms of \(\pi_x\) and taking their
masses gives maps \(T_j\) and weights \(\lambda_j\), after adding zero-weight
branches where necessary.  This proves \eqref{eq:finite-map-decomposition}
and \eqref{eq:finite-branch-characterization}.
\end{proof}

\begin{corollary}[Twist selects a unique graph coupling]
\label{prop:twist-linear-ot}
Under the assumptions of \cref{prop:finite-twist-linear-ot}, if \(c\) is
\(\alpha\)-almost-everywhere twisted on \(S_{f,g}\), then
\eqref{eq:linear-ot-value} has a unique optimizer
\((\Id,T)_\sharp\alpha\), characterized by
\begin{equation}
  \nabla f(x)=\nabla_xc(x,T(x))
  \qquad\text{for \(\alpha\)-almost every }x.
  \label{eq:generalized-brenier-characterization}
\end{equation}
\end{corollary}

\begin{proof}
The preceding proposition gives a graph optimizer.  If two optimal graph
couplings existed, their average would be optimal and the same argument would
force its conditional measures to be Dirac masses.  The two maps must
therefore agree \(\alpha\)-almost everywhere.
\end{proof}

Finite twist controls conditionals from the source side but generally gives
neither uniqueness nor an anti-graph structure.  The subtwist condition is a
different relaxation of twist that provides exactly these conclusions.

\begin{definition}[Subtwist]
\label{def:subtwist}
A cost \(c\), differentiable in its first variable on an open neighborhood
\(U\) of \(\supp\alpha\), satisfies \emph{subtwist} if, for every distinct
\(y_0,y_1\in\supp\beta\), the function
\begin{equation*}
  x\longmapsto c(x,y_0)-c(x,y_1)
\end{equation*}
has no critical point in \(U\), except possibly one global maximum and one
global minimum.
\end{definition}

\begin{theorem}[Subtwist selects a unique map/anti-map]
\label{thm:subtwist-linear-ot}
Assume compact supports, let \(c\) be continuous and \(C^1\) in its first
variable near the supports, and suppose that \(\alpha\) vanishes on every
\(C^1\) hypersurface.  If \(c\) satisfies subtwist, then the linear OT problem
\eqref{eq:linear-ot-value} has a unique optimizer, and this optimizer is a
map/anti-map in the sense of \cref{def:map-antimap}.
\end{theorem}

\begin{proof}[Proof mechanism]
This is the subtwist theorem of
\citet{ChiapporiMcCannNesheim2010,AhmadKimMcCann2011}.  On a dual contact set,
subtwist orders the possible contacts into a numbered limb system with at most
one graph limb and one anti-graph limb.  The hypothesis on \(\alpha\) excludes
mass on the hypersurfaces where the two limbs can switch.  A measure supported
on this acyclic two-limb system is uniquely determined by its marginals, which
gives both \eqref{eq:map-antimap-decomposition} and uniqueness.  Absolute
continuity \(\alpha\ll\mathrm{Leb}^n\), used throughout this note, is a
sufficient hypothesis.
\end{proof}

Combining this linear theorem with the universal freezing principle gives the
most general direct Brenier statement available for a quadratic GW energy.

\begin{corollary}[Abstract frozen-cost Brenier theorem]
\label{cor:abstract-frozen-cost-brenier}
Let \(Q\) be as in \eqref{eq:abstract-quadratic-gw}, assume
\(\alpha\ll\mathrm{Leb}^n\), and let \(\pi_\star\) minimize
\(\Ecal_Q\).  Suppose that its frozen cost \(c_{\pi_\star}\) is continuous
and \(C^1\) in its first variable near the compact supports.  If this cost is
\(m\)-twisted on a dual contact set for \(\alpha\)-almost every source
point, then \(\pi_\star\) is a \(k\)-map for some \(k\leq m\).  In
particular, one-twist implies that \(\pi_\star\) is a Monge coupling and is
the unique optimizer of its frozen linear problem.  No concavity of
\(\Ecal_Q\) is needed.
\end{corollary}

\begin{proof}
By \cref{prop:universal-frozen-gw}, \(\pi_\star\) minimizes linear OT for
\(c_{\pi_\star}\).  Apply
\cref{prop:finite-twist-linear-ot,prop:twist-linear-ot}.
\end{proof}

\subsection{The finite-branch power-distance theorem}

Freezing a GW optimizer now transfers the preceding linear result without
loss.  This is the appropriate Brenier statement when the twist map has more
than one inverse branch.

\begin{theorem}[Self-consistent finite-twist theorem]
\label{thm:rkhs-finite-twist}
Let \(0<p\leq2\), assume \(\alpha\ll\mathrm{Leb}^n\), and let
\(\pi_\star\) minimize \eqref{eq:power-gw-kantorovich}.  Set
\(A_\star=4C_{\pi_\star}\), let \((f_\star,g_\star)\) be a dual optimizer for
the frozen cost \(c_{A_\star}\), and suppose that this cost is \(C^1\) in
\(x\) near the supports.  If \(c_{A_\star}\) is
\(\alpha\)-almost-everywhere \(m\)-twisted on
\(S_{f_\star,g_\star}\), then \(\pi_\star\) is a \(k\)-map for some
\(k\leq m\), with active branches characterized by
\begin{equation}
  \nabla f_\star(x)=\nabla_xc_{A_\star}(x,T_j(x)).
  \label{eq:rkhs-finite-branch-characterization}
\end{equation}
If \(m=1\), the frozen linear optimizer is unique and \(\pi_\star\) is a
Monge coupling.
\end{theorem}

\begin{proof}
By \cref{cor:self-consistent-linear-ot}, \(\pi_\star\) solves linear OT for
\(c_{A_\star}\).  Apply \cref{prop:finite-twist-linear-ot}; the final claim is
\cref{prop:twist-linear-ot}.
\end{proof}

\begin{corollary}[RKHS-twist Brenier theorem for power-distance GW]
\label{thm:rkhs-twist-brenier}
Under the assumptions of \cref{thm:rkhs-finite-twist}, suppose that
\(c_{A_\star}\) is \(\alpha\)-almost-everywhere twisted.  Then
\begin{equation}
  \pi_\star=(\Id,T_\star)_\sharp\alpha,
  \label{eq:gw-monge-map}
\end{equation}
and it is the unique optimizer of the frozen problem.  Moreover,
\begin{equation}
  \nabla f_\star(x)
  =\nabla_xc_{A_\star}(x,T_\star(x))
  \quad\text{for \(\alpha\)-almost every }x.
  \label{eq:rkhs-brenier-characterization}
\end{equation}
\end{corollary}

This is a genuine Brenier-type conclusion, but the relevant inverse can be
multivalued: \eqref{eq:rkhs-finite-branch-characterization} selects the
branches of \(y\mapsto\nabla_xc_{A_\star}(x,y)\).  Only under one-twist is
there a single map, and that map need not be the gradient of an ordinary
convex function.

For \(1<p\leq2\), the twist map can be evaluated without constructing the
Hilbert features.  Define
\begin{equation}
  G_{\pi,x}(y)
  :=\int
  \norm{x-x'}^{p-2}(x-x')\norm{y-y'}^p
  \dd\pi(x',y').
  \label{eq:power-distance-twist-map}
\end{equation}
By \eqref{eq:gradient-direct-first-variation-cost}, the first-variation cost
has twist map \(-4pG_{\pi,x}\).

\begin{corollary}[Distance-only finite-twist criterion]
\label{cor:distance-only-twist}
Assume \(1<p\leq2\), \(\alpha\ll\mathrm{Leb}^n\), and compact supports.  Let
\(\pi_\star\) be GW-optimal.  If there is \(m<\infty\) such that, for
\(\alpha\)-almost every \(x\) and every \(q\in\R^n\),
\begin{equation}
  \card\{y\in\supp\beta:G_{\pi_\star,x}(y)=q\}\leq m,
  \label{eq:distance-m-twist}
\end{equation}
then \(\pi_\star\) is a \(k\)-map for some \(k\leq m\).  For \(m=1\), it is
a Monge coupling.  It is enough that \eqref{eq:distance-m-twist} hold after
intersecting each fiber with a dual contact set of the frozen problem.
\end{corollary}

\begin{proof}
For \(p>1\), the derivative of \(x\mapsto\norm{x-x'}^p\) is continuous,
including at \(x=x'\), and uniformly bounded on compact sets.  Thus
\(\ell_{\pi_\star}\) is continuously differentiable in \(x\), and
\eqref{eq:gradient-direct-first-variation-cost} identifies its gradient with
\(-4pG_{\pi_\star,x}\).  Since
\cref{prop:direct-first-variation-cost} makes \(\pi_\star\) optimal for this
frozen cost, \cref{prop:finite-twist-linear-ot} applies.
\end{proof}

\subsection{Quotient twist and transport inside the fibers}

Finite twist is a theorem about a \emph{given} optimal coupling.  A second
mechanism can prove existence of \emph{some} Monge optimizer even when every
twist fiber is infinite.  It separates the variables seen by the cost from
the directions along which the cost is indifferent.  One first transports
the visible quotient variables and then uses the invisible source fibers to
reproduce the target conditional laws.  The measurable lifting result below
is the abstract mechanism developed by
\citet[Theorem~2.1]{DumontLacombeVialard2025}.

\begin{theorem}[Quotient-and-fiber lifting for linear OT]
\label{thm:quotient-fiber-linear-ot}
Let \(\mathcal X,\mathcal Y,\mathcal U,\mathcal V\) be standard Borel
spaces, let
\(
  \phi:\mathcal X\to\mathcal U
\)
and
\(
  \psi:\mathcal Y\to\mathcal V
\)
be measurable, and suppose that
\begin{equation}
  c(x,y)=\widetilde c\bigl(\phi(x),\psi(y)\bigr).
  \label{eq:quotient-cost-factorization}
\end{equation}
Set \(\bar\alpha=\phi_\sharp\alpha\) and
\(\bar\beta=\psi_\sharp\beta\).  Assume that quotient OT for
\(\widetilde c\) between \(\bar\alpha\) and \(\bar\beta\) has an optimal
map \(t:\mathcal U\to\mathcal V\), and let
\begin{equation}
  \alpha=\int\alpha_u\dd\bar\alpha(u),
  \qquad
  \beta=\int\beta_v\dd\bar\beta(v)
  \label{eq:quotient-disintegrations}
\end{equation}
be disintegrations along \(\phi\) and \(\psi\).  If \(\alpha_u\) is
nonatomic for \(\bar\alpha\)-almost every \(u\), or more generally if there
is a jointly measurable family of maps satisfying
\begin{equation}
  (T_u)_\sharp\alpha_u=\beta_{t(u)}
  \quad\text{for \(\bar\alpha\)-almost every \(u\)},
  \label{eq:fiberwise-transport-maps}
\end{equation}
then there exists a
measurable map \(T:\mathcal X\to\mathcal Y\) such that
\begin{equation}
  T_\sharp\alpha=\beta,
  \qquad
  \psi\bigl(T(x)\bigr)=t\bigl(\phi(x)\bigr)
  \quad\text{for \(\alpha\)-almost every \(x\)},
  \label{eq:quotient-map-lift}
\end{equation}
and \((\Id,T)_\sharp\alpha\) is optimal for the original cost \(c\).
\end{theorem}

\begin{proof}
When the fiber maps are not given directly, the measurable lifting theorem of
\citet[Theorem~2.1]{DumontLacombeVialard2025} supplies them.  Indeed, every
standard Borel space is Borel-isomorphic to a Borel subset of \([0,1]\), so
that result applies to the probability kernels
\(u\mapsto\alpha_u\) and \(u\mapsto\beta_{t(u)}\).  Atomlessness of the first
kernel is exactly what permits it to generate the second one by a jointly
measurable family of deterministic maps.  Set
\(T(x)=T_{\phi(x)}(x)\).  Then
\begin{align*}
  T_\sharp\alpha
  &=\int (T_u)_\sharp\alpha_u\dd\bar\alpha(u)
    =\int\beta_{t(u)}\dd\bar\alpha(u)\\
  &=\int\beta_v\dd\bar\beta(v)=\beta,
\end{align*}
because \(t_\sharp\bar\alpha=\bar\beta\).  This also proves the projection
identity in \eqref{eq:quotient-map-lift}.

For any \(\pi\in\Piopt(\alpha,\beta)\), the pushforward
\((\phi,\psi)_\sharp\pi\) belongs to
\(\Piopt(\bar\alpha,\bar\beta)\), and
\begin{equation*}
  \int c\dd\pi
  =\int\widetilde c\dd\bigl((\phi,\psi)_\sharp\pi\bigr).
\end{equation*}
The graph coupling constructed above projects to
\((\Id,t)_\sharp\bar\alpha\), so it attains the quotient optimum and is
therefore optimal for \(c\).
\end{proof}

Atomlessness is therefore a clean sufficient condition, not a logically
necessary one.  For example, no atomlessness is needed when
\(\beta_{t(u)}\) is a Dirac mass.  Conversely, if both conditional laws have
incompatible atoms, quotient optimality need not lift to a map.  This fiber
compatibility is exactly what distinguishes a quotient argument from a twist
argument.

The replacement principle in \cref{prop:universal-frozen-gw} now turns this
linear lifting theorem into an abstract GW Brenier theorem.

\begin{theorem}[Quotient-and-fiber Brenier principle for GW]
\label{thm:quotient-fiber-gw}
Let \(Q\) satisfy the conditional concavity condition
\eqref{eq:conditional-concavity-gw-kernel}, let \(\pi_\star\) minimize
\eqref{eq:abstract-quadratic-gw}, and let \(c_{\pi_\star}\) be its frozen
cost \eqref{eq:abstract-frozen-gw-cost}.  Suppose that this cost admits a
factorization \eqref{eq:quotient-cost-factorization} for which the hypotheses
of \cref{thm:quotient-fiber-linear-ot} hold.  Then \(\Ecal_Q\) has a Monge
minimizer \((\Id,T)_\sharp\alpha\).

If \(\mathcal U,\mathcal V\) are compact subsets of \(\R^r\), the quotient
cost is continuous, continuously differentiable in its first variable near
the supports, and twisted there, and
\(\bar\alpha\ll\mathrm{Leb}^r\), then the required quotient map exists and
is unique.
\end{theorem}

\begin{proof}
By \cref{prop:universal-frozen-gw}, \(\pi_\star\) minimizes linear OT for
\(c_{\pi_\star}\).  The quotient-and-fiber theorem constructs a graph
coupling \(\widehat\pi=(\Id,T)_\sharp\alpha\) that minimizes the same frozen
problem.  Conditional concavity and the replacement part of
\cref{prop:universal-frozen-gw} imply that \(\widehat\pi\) is also a global
minimizer of \(\Ecal_Q\).  The final assertion is the usual twisted-cost
Monge theorem applied on the quotient spaces.
\end{proof}

The concavity hypothesis is essential for this \emph{replacement} argument.
Without it, the original optimizer \(\pi_\star\) still solves its frozen
linear problem, but another optimizer of that linear problem need not remain
GW-optimal.  Thus one-twist of the full frozen cost proves that
\(\pi_\star\) itself is a map without concavity, whereas quotient lifting
proves existence of a possibly different map by using concavity.
For a particular lifted coupling \(\widehat\pi\), global conditional
concavity can be weakened to the exact directional condition
\begin{equation}
  \iint Q(z,z')
  \dd(\widehat\pi-\pi_\star)(z)
  \dd(\widehat\pi-\pi_\star)(z')=0,
  \label{eq:directional-quotient-replacement}
\end{equation}
by \eqref{eq:frozen-replacement-defect}.  Condition
\eqref{eq:conditional-concavity-gw-kernel} is valuable because it enforces
this identity automatically for every frozen lift.
Nothing in \cref{thm:quotient-fiber-gw} requires a linear or finite-rank
factorization: any finite-dimensional sufficient statistic
\((\phi,\psi)\) for the frozen cost is admissible.  Finite rank is treated
next because it supplies such statistics canonically for the learned RKHS
cost.

\subsection{Finite-rank quotient lifting and feature slices}

Condition \eqref{eq:distance-m-twist} is exact but still depends on the
unknown optimizer.  Finite rank gives two complementary routes.  Its
quotient variables can produce a Monge optimizer by the preceding lifting
theorem, while their affine slices control the number of branches of a
prescribed optimizer.

Suppose first that the learned operator has finite rank \(R\), with singular
value decomposition
\begin{equation}
  A=\sum_{j=1}^R\sigma_j e_j\otimes f_j,
  \qquad
  \varphi_j(x):=\ip{u(x)}{f_j},
  \qquad
  \psi_j(y):=\ip{v(y)}{e_j}.
  \label{eq:finite-rank-learned-operator}
\end{equation}
Define the visible source and target statistics
\begin{align}
  \Phi_A(x)&:=\bigl(a(x),\varphi_1(x),\ldots,\varphi_R(x)\bigr),
  &
  \Psi_A(y)&:=\bigl(b(y),\psi_1(y),\ldots,\psi_R(y)\bigr),
  \label{eq:finite-rank-quotient-features}\\
  D_A&:=\operatorname{diag}(1,\sigma_1,\ldots,\sigma_R).
  \label{eq:finite-rank-diagonal-operator}
\end{align}
Then the learned cost factors through an \((R+1)\)-dimensional bilinear
problem:
\begin{equation}
  c_A(x,y)=-4\ip{D_A\Phi_A(x)}{\Psi_A(y)}.
  \label{eq:finite-rank-bilinear-quotient-cost}
\end{equation}

\begin{theorem}[Finite-rank quotient Brenier theorem]
\label{thm:finite-rank-quotient-brenier}
Let \(0<p\leq2\), let \(\pi_\star\) be GW-optimal, and suppose that
\(A_\star=4C_{\pi_\star}\) has finite rank \(R\).  Assume that
\begin{equation}
  (\Phi_{A_\star})_\sharp\alpha\ll\mathrm{Leb}^{R+1}
  \label{eq:finite-rank-quotient-density}
\end{equation}
and that the conditional measures in a disintegration of \(\alpha\) along
\(\Phi_{A_\star}\) are nonatomic almost everywhere.  Then there exists a
GW-optimal map \(T_\star:\mathcal X\to\mathcal Y\).  It may be chosen so
that
\begin{equation}
  \Psi_{A_\star}\bigl(T_\star(x)\bigr)
  =t_\star\bigl(\Phi_{A_\star}(x)\bigr),
  \label{eq:finite-rank-lifted-map}
\end{equation}
where \(t_\star\) is the unique quotient OT map for the bilinear cost in
\eqref{eq:finite-rank-bilinear-quotient-cost}.  Moreover, the lifted
optimizer has the same cross-covariance operator as \(\pi_\star\).
\end{theorem}

\begin{proof}
All singular values in \eqref{eq:finite-rank-diagonal-operator} are positive,
so \((D_{A_\star})_\sharp(\Phi_{A_\star})_\sharp\alpha\) has a density.
Moreover,
\begin{equation*}
  -4\ip{D_{A_\star}s}{z}
  =2\norm{D_{A_\star}s-z}^2
   -2\norm{D_{A_\star}s}^2-2\norm{z}^2.
\end{equation*}
The last two terms are fixed by the quotient marginals, so the classical
Brenier theorem gives a unique quotient map \(t_\star\); equivalently,
\(t_\star(s)=\nabla h(D_{A_\star}s)\) for a convex function \(h\).
The factorization
\eqref{eq:finite-rank-bilinear-quotient-cost} and
the concavity established in \eqref{eq:gw-feature-decomposition} allow
\cref{thm:quotient-fiber-gw} to give a GW-optimal lift.

It remains to verify the last assertion.  Put
\(
  \eta=(\Id,T_\star)_\sharp\alpha-\pi_\star.
\)
Extend \eqref{eq:hilbert-cross-covariance} linearly to signed measures, so
that \(C_\eta=C_{(\Id,T_\star)_\sharp\alpha}-C_{\pi_\star}\).
The two couplings tie for the frozen linear objective and for the GW
objective.  By \eqref{eq:gw-feature-decomposition}, the quadratic coefficient
of \(t\mapsto\Ecal_p(\pi_\star+t\eta)\) is
\(-8\norm{C_\eta}_{\HS}^2\).  The segment expansion
\eqref{eq:quadratic-gw-segment-expansion} consequently forces
\(C_\eta=0\), proving
\(C_{(\Id,T_\star)_\sharp\alpha}=C_{\pi_\star}\).
\end{proof}

The two measure-theoretic hypotheses have a useful differential sufficient
condition.

\begin{corollary}[A coarea criterion]
\label{cor:finite-rank-coarea-brenier}
In addition to the assumptions of
\cref{thm:finite-rank-quotient-brenier}, suppose that
\(\mathcal X\subset\R^n\), \(\alpha\ll\mathrm{Leb}^n\),
\(\Phi_{A_\star}\) is \(C^1\) near \(\supp\alpha\), and
\begin{equation}
  n>R+1,
  \qquad
  \rank D\Phi_{A_\star}(x)=R+1
  \quad\text{for \(\alpha\)-almost every \(x\)}.
  \label{eq:finite-rank-coarea-condition}
\end{equation}
Then a GW-optimal Monge map exists.
\end{corollary}

\begin{proof}
Write \(\dd\alpha(x)=\rho_\alpha(x)\dd x\) and
\(\Phi=\Phi_{A_\star}\).  Since \(\Phi\) is \(C^1\) near the compact source
support, it is locally Lipschitz there.  Its \((R+1)\)-dimensional Jacobian
\(J_{R+1}\Phi\) is positive at almost every source point by
\eqref{eq:finite-rank-coarea-condition}.  The coarea formula
\citep{EvansGariepy2015} gives, for bounded measurable \(h\),
\begin{equation}
  \int h(x)\rho_\alpha(x)\dd x
  =\int_{\R^{R+1}}
    \int_{\Phi^{-1}(s)}
      h(x)\frac{\rho_\alpha(x)}{J_{R+1}\Phi(x)}
      \dd\mathcal H^{n-R-1}(x)\dd s,
  \label{eq:finite-rank-coarea-disintegration}
\end{equation}
where the \(\alpha\)-null critical set may be discarded.  Taking \(h\) that
depends only on \(s=\Phi(x)\) shows that \(\Phi_\sharp\alpha\) has a density.
After normalizing the inner measures in
\eqref{eq:finite-rank-coarea-disintegration}, one obtains a disintegration of
\(\alpha\) along \(\Phi\), absolutely continuous with respect to
\(\mathcal H^{n-R-1}\) on almost every regular fiber.  Since
\(n-R-1\geq1\), these conditional measures are nonatomic.  Apply
\cref{thm:finite-rank-quotient-brenier}.
\end{proof}

The coarea criterion is not a finite-twist condition: its regular fibers have
positive dimension.  Those fibers are precisely the reservoir of randomness
used to match the target conditional laws deterministically.

Whenever \(a\) and the \(\varphi_j\) are differentiable, define
\begin{align}
  Z_A(y)&:=\Psi_A(y)\in\R^{R+1},
  \label{eq:augmented-target-feature}\\
  B_{A,x}(s_0,\ldots,s_R)
    &:=s_0\nabla a(x)+\sum_{j=1}^R
      s_j\sigma_j\nabla\varphi_j(x)
      \in\R^n.
  \label{eq:feature-projection-operator}
\end{align}

\begin{proposition}[Finite-rank feature-slice criterion]
\label{prop:finite-rank-feature-slices}
Under the preceding assumptions, the learned twist map factors as
\begin{equation}
  \nabla_xc_A(x,y)=-4B_{A,x}Z_A(y).
  \label{eq:finite-rank-twist-factorization}
\end{equation}
Consequently, \(c_A\) is \(m\)-twisted at \(x\) whenever every affine fiber
of \(B_{A,x}\) intersects \(Z_A(\mathcal Y)\) in at most \(m\) points after
pulling the intersection back through \(Z_A\).  If, in addition, \(Z_A\) has
fibers of cardinality at most \(s\), the complex projective closure of
\(Z_A(\mathcal Y)\) is an algebraic variety of degree \(D\), and these linear
sections are zero dimensional, then one may take \(m\leq sD\).
\end{proposition}

\begin{proof}
The singular value decomposition gives
\(
 \ip{Au(x)}{v(y)}=\sum_j\sigma_j\varphi_j(x)\psi_j(y).
\)
Differentiating \eqref{eq:learned-cost} in \(x\) proves
\eqref{eq:finite-rank-twist-factorization}.  The first conclusion is then
just a fiber identity.  The algebraic conclusion is the definition of degree:
a zero-dimensional linear section of a projective variety contains at most
\(D\) complex points when multiplicities are counted.  Each image point has
at most \(s\) preimages under \(Z_A\), which gives \(sD\).
\end{proof}

For \(p=2\) and full row rank, the target feature set is, after an invertible
linear change of coordinates, the paraboloid
\begin{equation}
  Z_A(y)=\bigl(\norm{y-m_\beta}^2,y-m_\beta\bigr),
  \label{eq:p2-augmented-feature-paraboloid}
\end{equation}
whose projective closure has degree two.  The finite-rank criterion therefore
predicts at most two branches; \cref{lem:p2-rank-geometry} below verifies the
zero-dimensionality condition and locates exactly where the second branch can
occur.  For \(0<p<2\), the optimal Hilbert--Schmidt operator need not have
finite rank, and even a finite-rank truncation does not preserve exact
optimality.  Thus this argument identifies a genuine rigidity route, but not
an automatic finite-map theorem.

For quadratic dissimilarities, the coarea criterion recovers the low-rank
Monge construction and sharpens it when the ambient source has additional
dimensions.

\begin{corollary}[Quadratic finite-rank quotient regimes]
\label{cor:p2-finite-rank-quotient}
Let \(p=2\), let \(n\geq d\), assume
\(\alpha\ll\mathrm{Leb}^n\), and let \(\pi_\star\) be GW-optimal with
\begin{equation}
  h:=\rank S_{\pi_\star}.
  \label{eq:p2-optimal-rank-h}
\end{equation}
If \(n>h+1\), then there exists a GW-optimal Monge map with the same
cross-covariance as \(\pi_\star\).  Consequently, this conclusion holds
whenever either
\begin{equation}
  h\leq d-2,
  \qquad\text{or}\qquad
  h=d-1\ \text{ and }\ n>d.
  \label{eq:p2-quotient-rank-regimes}
\end{equation}
\end{corollary}

\begin{proof}
Up to orthogonal coordinates, the source quotient statistics are
\begin{equation}
  \Phi_{A_\star}(x)
  =\bigl(\norm{x-m_\alpha}^2,
          \ip{x-m_\alpha}{f_1},\ldots,
          \ip{x-m_\alpha}{f_h}\bigr),
  \label{eq:p2-source-quotient-map}
\end{equation}
where \(f_1,\ldots,f_h\) are orthonormal.  Its differential has rank
\(h+1\) unless
\(
  x-m_\alpha\in\operatorname{span}\{f_1,\ldots,f_h\}.
\)
This exceptional affine subspace is \(\alpha\)-null.  If \(n>h+1\),
\cref{cor:finite-rank-coarea-brenier} applies.  Finally,
\(h\leq d-2\) and \(n\geq d\) imply \(n>h+1\), while
\(h=d-1\) and \(n>d\) give the same strict inequality.
\end{proof}

For \(h\leq d-2\), this is the quotient mechanism behind the map theorem of
\citet{DumontLacombeVialard2025}.  When \(h=d-1\), equal dimensions leave
zero-dimensional source and target fibers whose atomic weights need not be
compatible; the general theorem then gives only two branches.  The strict
dimension gap turns the source fibers into positive-dimensional spheres, so
they become nonatomic and a single-map lift exists.  Full row rank
\(h=d\) is treated more efficiently by the direct twist argument in
\cref{thm:p2-dimension-gap-monge}, including the borderline case \(n=d+1\).

The borderline zero-dimensional fibers still yield a map when the target
feature separates its support.  This gives a geometric refinement of the
rank-only statement.

\begin{proposition}[Injective radial-linear target features]
\label{prop:p2-injective-target-feature}
Let \(p=2\), assume compact supports and
\(\alpha\ll\mathrm{Leb}^n\), and let \(\pi_\star\) be GW-optimal.  Put
\begin{equation*}
  h:=\rank S_{\pi_\star},
  \qquad
  E_\star:=\Ran(S_{\pi_\star})\subset\R^d,
\end{equation*}
and assume \(h<n\).  Define
\begin{equation}
  \Theta_{E_\star}(y)
  :=\left(\norm{y-m_\beta}^2,
           P_{E_\star}(y-m_\beta)\right)
  \in\R\times E_\star,
  \label{eq:p2-target-radial-linear-feature}
\end{equation}
where \(P_{E_\star}\) is orthogonal projection.  If
\(\Theta_{E_\star}\) is injective on \(\supp\beta\), then
\(\pi_\star\) itself is induced by a map and is the unique optimizer of its
frozen linear OT problem.
\end{proposition}

\begin{proof}
Let \(F_\star:=\Ran(S_{\pi_\star}^\top)\subset\R^n\).  In orthonormal bases
of \(F_\star\) and \(E_\star\), the source and target quotient variables in
\eqref{eq:finite-rank-quotient-features} are the basis-free maps
\begin{equation*}
  \Theta_{F_\star}(x)
  =\left(\norm{x-m_\alpha}^2,P_{F_\star}(x-m_\alpha)\right),
  \qquad
  \Theta_{E_\star}(y).
\end{equation*}
The differential of \(\Theta_{F_\star}\) has rank \(h+1\) unless
\(x-m_\alpha\in F_\star\).  Because \(h<n\), this exceptional affine
subspace is \(\alpha\)-null.  The coarea formula therefore gives
\begin{equation*}
  (\Theta_{F_\star})_\sharp\alpha
  \ll\mathrm{Leb}^{h+1}.
\end{equation*}

The frozen cost factors as the bilinear quotient cost
\eqref{eq:finite-rank-bilinear-quotient-cost}.  After the invertible diagonal
change \(D_{A_\star}\), Brenier's theorem makes its quotient optimizer unique
and induced by a map \(t_\star\).  Every quotient coupling lifts to the
original marginals by coupling the corresponding conditional measures, so
the projection of every frozen optimizer must solve this quotient problem.
In particular, \(\pi_\star\) satisfies
\begin{equation*}
  \Theta_{E_\star}(y)
  =t_\star\bigl(\Theta_{F_\star}(x)\bigr)
  \quad\text{for \(\pi_\star\)-almost every \((x,y)\)}.
\end{equation*}
Since \(\Theta_{E_\star}\) is a continuous injection from the compact set
\(\supp\beta\), its inverse on its image is measurable.  Thus the last
display determines
\begin{equation*}
  y=T_\star(x)
  :=\Theta_{E_\star}^{-1}\!\left(
      t_\star(\Theta_{F_\star}(x))\right)
\end{equation*}
for \(\pi_\star\)-almost every \((x,y)\), proving that \(\pi_\star\) is a
graph.  The same argument applies to every optimizer of the frozen problem,
and the unique quotient map and injective target feature force the same
graph; hence frozen optimality is unique.
\end{proof}

The fibers in \eqref{eq:p2-target-radial-linear-feature} are easy to see:
after fixing the projection onto \(E_\star\), they are spheres in
\(E_\star^\perp\).  Thus the hypothesis says that \(\supp\beta\) meets each
such sphere in at most one point.  In the equal-dimensional borderline case
\(n=d\), \(h=d-1\), each nondegenerate fiber consists of two points related by
the affine reflection
\begin{equation*}
  \mathfrak r_{E_\star}(y)
  :=m_\beta+2P_{E_\star}(y-m_\beta)-(y-m_\beta).
\end{equation*}
Consequently, the concrete condition
\begin{equation}
  y,\mathfrak r_{E_\star}(y)\in\supp\beta
  \quad\Longrightarrow\quad y\in m_\beta+E_\star
  \label{eq:p2-no-reflected-target-pairs}
\end{equation}
upgrades the general two-map conclusion to a Monge optimizer.  This identifies
the precise target geometry responsible for the second branch.  The criterion
is a direct refinement obtained from the quotient formalism above; it does not
replace the unconditional two-map theorem of
\citet{DumontLacombeVialard2025}.

\subsection{Differential, generic, and definable finite-twist criteria}

\begin{proposition}[Mixed-Hessian immersion gives finite twist]
\label{prop:mixed-hessian-finite-twist}
Let \(\mathcal X\subset\R^n\) be compact and let \(\mathcal Y\) be a compact
\(d\)-dimensional \(C^1\) manifold, possibly with boundary.  Suppose
\(F(x,y):=\nabla_xc(x,y)\) is continuous and continuously differentiable in
\(y\), and
\begin{equation}
  \rank D_yF(x,y)=d
  \qquad\text{for every }(x,y)\in\mathcal X\times\mathcal Y.
  \label{eq:mixed-hessian-full-column-rank}
\end{equation}
Then there is an integer \(m<\infty\) such that every fiber of
\(y\mapsto F(x,y)\) has at most \(m\) points, uniformly in \(x\).  Hence \(c\)
is globally \(m\)-twisted.
\end{proposition}

\begin{proof}
At each \((x_0,y_0)\), choose \(d\) output coordinates for which the
corresponding \(d\times d\) minor of \(D_yF(x_0,y_0)\) is invertible.  The
inverse-function theorem with parameters gives product neighborhoods
\(U\times V\) of \((x_0,y_0)\) on which, for every fixed \(x\in U\), the
restriction \(F_x|_V\) is injective.  Compactness provides a finite cover by
such product neighborhoods.  A fiber meets each member of this cover in at
most one point, so its cardinality is bounded by the number of members in the
cover.
\end{proof}

For the power-distance cost, the mixed Hessian has a particularly transparent
form.  With vectors understood as columns, set
\begin{equation}
  M_{\pi,p}(x,y)
  :=\int
  \norm{x-x'}^{p-2}\norm{y-y'}^{p-2}
  (x-x')(y-y')^\top\dd\pi(x',y')
  \in\R^{n\times d},
  \label{eq:power-cross-hessian}
\end{equation}
using the continuous zero extension of
\(z\mapsto\norm{z}^{p-2}z\) at \(z=0\).

\begin{corollary}[A cross-Hessian criterion for finite-branch GW]
\label{cor:power-cross-hessian-finite-map}
Let \(1<p\leq2\), let \(\alpha\ll\mathrm{Leb}^n\), and suppose the compact
target support lies in a compact \(d\)-dimensional \(C^1\) manifold.  If a
GW optimizer \(\pi_\star\) satisfies
\begin{equation}
  \rank M_{\pi_\star,p}(x,y)=d
  \qquad\text{on }\supp\alpha\times\supp\beta,
  \label{eq:power-cross-hessian-rank-condition}
\end{equation}
then \(\pi_\star\) is a \(k\)-map for some finite \(k\).
\end{corollary}

\begin{proof}
Differentiating \eqref{eq:power-distance-twist-map} in \(y\) gives
\begin{equation}
  D_yG_{\pi,x}(y)=pM_{\pi,p}(x,y),
  \qquad
  D_y\nabla_x\ell_\pi(x,y)=-4p^2M_{\pi,p}(x,y).
  \label{eq:power-mixed-hessian-identity}
\end{equation}
The integrand is continuous for \(p>1\), so differentiation under the
integral is valid on compact sets.  Apply
\cref{prop:mixed-hessian-finite-twist,cor:distance-only-twist}.
\end{proof}

Condition \eqref{eq:power-cross-hessian-rank-condition} is strong: it rules
out folds of the projected feature geometry.  A weaker, generic mechanism
counts how many sheets a smooth map can have.  For fixed \(x\), write
\(F_x(y)=\nabla_xc(x,y)\), and for pairwise distinct
\((y_1,\ldots,y_r)\) define the coincidence map
\begin{equation}
  \Delta_rF_x(y_1,\ldots,y_r)
  :=\bigl(F_x(y_2)-F_x(y_1),\ldots,
          F_x(y_r)-F_x(y_1)\bigr)
  \in(\R^n)^{r-1}.
  \label{eq:multijet-coincidence-map}
\end{equation}

\begin{proposition}[General-position branch count]
\label{prop:general-position-branch-count}
Assume the hypotheses of \cref{prop:finite-twist-linear-ot}, let \(n>d\),
let \(\mathcal Y\) be a compact smooth \(d\)-manifold, and suppose
\(F_x(y)=\nabla_xc(x,y)\) is smooth in \(y\).  Set
\begin{equation}
  m_{n,d}:=\left\lfloor\frac{n}{n-d}\right\rfloor.
  \label{eq:generic-branch-number}
\end{equation}
If, for \(\alpha\)-almost every \(x\), the map
\(\Delta_{m_{n,d}+1}F_x\) is transverse to the origin on the set of pairwise
distinct tuples, then every fiber of \(F_x\) has at most \(m_{n,d}\) points.
Consequently, a frozen optimal plan for \(c\) is an
\(m_{n,d}\)-map.  In particular, general position predicts one branch when
\(n>2d\), at most two when \(n=2d\), and at most three when \((n,d)=(3,2)\).
\end{proposition}

\begin{proof}
The domain of \(\Delta_rF_x\) has dimension \(rd\), whereas its target has
dimension \((r-1)n\).  For \(r=m_{n,d}+1\),
\[
  rd-(r-1)n=d-m_{n,d}(n-d)<0.
\]
A map from a lower-dimensional manifold cannot be transverse to the origin at
a preimage of the origin.  Hence no \(r\) distinct points have the same
image, which is exactly the asserted fiber bound.  The finite-map conclusion
follows from \cref{prop:finite-twist-linear-ot}.
\end{proof}

The restriction \(n>d\) is structural for this argument.  If \(n<d\), a
regular fiber of a smooth map \(\mathcal Y^d\to\R^n\) has dimension
\(d-n\), so finite twist is nongeneric.  If \(n=d\), regular fibers are
discrete but their number is a topological or algebraic degree not controlled
by dimension alone.

Multijet transversality makes the hypothesis of
\cref{prop:general-position-branch-count} generic for unconstrained smooth
maps \citep{GolubitskyGuillemin1973}.  It is not automatically generic inside
the much smaller family of self-consistent GW costs.  Proving transversality
with respect to perturbations of \((\alpha,\beta)\), while retaining the
fixed-point relation \(A_\star=4C_{\pi_\star}\), is a concrete route to a
generic finite-map theorem.  For \(1<p<2\), smoothness of
\(G_{\pi_\star,x}\) itself requires regularity of the vector-valued measures
obtained by weighting \(\pi_\star\); smoothness of \(\beta\) alone does not
automatically provide it.

Analyticity offers a different rigidity mechanism, but only after
positive-dimensional fibers have been excluded.

\begin{proposition}[Definable analytic rigidity]
\label{prop:definable-finite-twist}
Let \(\mathcal X\), \(\mathcal Y\), and
\(F:\mathcal X\times\mathcal Y\to\R^n\) be definable in a fixed o-minimal
structure, with \(\mathcal Y\) compact.  Define
\begin{equation}
  E:=\{x\in\mathcal X:\text{some fiber }
       \{y\in\mathcal Y:F(x,y)=q\}\text{ is infinite}\}.
  \label{eq:infinite-fiber-exceptional-set}
\end{equation}
Then \(E\) is definable and there exists \(m<\infty\) such that
\begin{equation}
  \card\{y\in\mathcal Y:F(x,y)=q\}\leq m
  \qquad(x\notin E,\ q\in\R^n).
  \label{eq:definable-uniform-finiteness}
\end{equation}
In particular, if \(F(x,y)=\nabla_xc(x,y)\) and \(\alpha(E)=0\), then \(c\)
is \(\alpha\)-almost-everywhere \(m\)-twisted.  This applies to a
real-analytic cost on neighborhoods of compact semianalytic supports,
provided no gradient fiber has positive dimension outside an
\(\alpha\)-null set.
\end{proposition}

\begin{proof}
In an o-minimal structure, an infinite definable fiber has positive
dimension, and the locus where the fiber dimension is positive is definable.
On the complement of this locus all fibers are finite.  The uniform
finiteness theorem for definable families then supplies a common cardinality
bound \citep{vandenDries1998}.  Restricted real-analytic functions on compact
semianalytic sets are globally subanalytic, hence definable in
\(\R_{\mathrm{an}}\).
\end{proof}

There are two important qualifications.  First, analyticity alone neither
excludes infinite fibers nor fixes their number.  On \([-1,1]^2\), the
analytic costs
\(
 c_N(x,y)=xT_N(y)
\)
built from Chebyshev polynomials satisfy
\(\partial_xc_N=T_N(y)\), whose typical fibers have \(N\) points; a cost
independent of \(y\) has infinite fibers.  Second, when \(0<p<2\), the kernel
\(\norm{y-y'}^p\) is not analytic on the diagonal, and integration against an
arbitrary optimal coupling does not preserve definability.  Thus
\cref{prop:definable-finite-twist} is a conditional rigidity tool, not an
automatic consequence of the Schoenberg representation.

In algebraic models the branch number can sometimes be made explicit.  If
the components of \(F_x(y)\) are polynomials of degree at most \(r\) and every
complexified fiber is zero dimensional, B\'ezout's theorem gives the crude
uniform bound \(r^d\).  For \(p=2\), the frozen twist map is quadratic in
\(y\), so this argument gives \(2^d\); the rank-one structure analyzed next
improves it sharply to two.

\section{The finite-dimensional case \texorpdfstring{\(p=2\)}{p=2}}
\label{sec:p2}

At \(p=2\), the abstract feature objects become ordinary centered moments.
This permits an exact analysis of the obstruction.

\subsection{The learned finite-dimensional cost}

Write
\begin{equation}
  \widetilde x=x-m_\alpha,
  \qquad
  \widetilde y=y-m_\beta.
  \label{eq:centered-euclidean-variables}
\end{equation}
Then
\[
  u(x)=\widetilde x,
  \quad v(y)=\widetilde y,
  \quad a(x)=\norm{\widetilde x}^2,
  \quad b(y)=\norm{\widetilde y}^2,
  \quad C_\pi=S_\pi.
\]
The general cross-Hessian \eqref{eq:power-cross-hessian} reduces to
\begin{equation}
  M_{\pi,2}(x,y)
  =\widetilde x\widetilde y^\top+S_\pi^\top.
  \label{eq:p2-cross-hessian}
\end{equation}
The operator \(A\) is a matrix in \(\R^{d\times n}\), the RKHS penalty is
the Frobenius norm, and
\begin{equation}
  c_A(x,y)
  =-4\norm{\widetilde x}^2\norm{\widetilde y}^2
   -4\ip{A\widetilde x}{\widetilde y},
  \qquad
  A_\star=4S_{\pi_\star}.
  \label{eq:p2-learned-cost}
\end{equation}
The first term is the crucial difference from inner-product GW.  Without it,
full row rank of \(A\) immediately gives twist.  With it, the gradient is
\begin{equation}
  \nabla_xc_A(x,y)
  =-8\widetilde x\norm{\widetilde y}^2
   -4A^\top\widetilde y.
  \label{eq:p2-cost-gradient}
\end{equation}

\subsection{Rotational symmetry gives a global radial Brenier theorem}

Equal-dimensional rotational symmetry removes the two-branch obstruction
globally, without any condition on an unknown optimal cross-covariance.  The
result is due to \citet[Theorem~9.21]{Sturm2023Space}; the formulation below
makes its two rearrangement mechanisms explicit in the notation of this note.

In this subsection only, compact support is replaced by the natural moment
assumption
\begin{equation}
  \int\norm{x}^4\dd\alpha(x)+\int\norm{y}^4\dd\beta(y)<\infty.
  \label{eq:p2-radial-fourth-moments}
\end{equation}
Assume that \(\alpha,\beta\in\Pcal(\R^d)\) satisfy
\eqref{eq:p2-radial-fourth-moments} and are invariant under every orthogonal
transformation around their respective barycenters.  Define their radial laws
by
\begin{equation}
  \nu_\alpha
  :=\bigl(x\mapsto\norm{x-m_\alpha}\bigr)_\sharp\alpha,
  \qquad
  \nu_\beta
  :=\bigl(y\mapsto\norm{y-m_\beta}\bigr)_\sharp\beta.
  \label{eq:p2-radial-laws}
\end{equation}
Let \(F_\alpha^{\mathrm{rad}}\) be the cumulative function of
\(\nu_\alpha\), let \(Q_\beta^{\mathrm{rad}}\) be the quantile function of
\(\nu_\beta\), and put
\begin{equation}
  \tau(r):=
  Q_\beta^{\mathrm{rad}}\bigl(F_\alpha^{\mathrm{rad}}(r)\bigr),
  \qquad
  \mathsf M_k:=\int_0^\infty r^k\tau(r)^k\dd\nu_\alpha(r),
  \quad k\in\{1,2\}.
  \label{eq:p2-radial-rearrangement}
\end{equation}
When \(\alpha\ll\mathrm{Leb}^d\), its radial law is atomless and
\(\tau_\sharp\nu_\alpha=\nu_\beta\).

\begin{lemma}[Sturm's simultaneous radial and angular maximization]
\label{lem:sturm-polar-maximization}
Assume that \(\alpha,\beta\ll\mathrm{Leb}^d\) satisfy
\eqref{eq:p2-radial-fourth-moments} and are rotationally invariant around
their barycenters.  For every \(\pi\in\Piopt(\alpha,\beta)\),
\begin{equation}
  \int\norm{x-m_\alpha}^2\norm{y-m_\beta}^2\dd\pi(x,y)
  \leq\mathsf M_2,
  \qquad
  \norm{S_\pi}_{\mathrm F}^2\leq\frac{\mathsf M_1^2}{d}.
  \label{eq:p2-sturm-two-bounds}
\end{equation}
Simultaneous equality holds if and only if, for some \(O\in\mathrm O(d)\),
\begin{equation}
  \pi=(\Id,T_O)_\sharp\alpha,
  \qquad
  T_O(m_\alpha+r\omega)
  :=m_\beta+\tau(r)O\omega
  \quad
  \text{for }r>0,\ \omega\in\mathbb S^{d-1}.
  \label{eq:p2-sturm-radial-map}
\end{equation}
The value of \(T_O\) at \(m_\alpha\) is immaterial.
\end{lemma}

\begin{proof}
Center both laws and use polar coordinates
\(x=r\omega\), \(y=s\zeta\), and denote the resulting law of
\((r,\omega,s,\zeta)\) by \(\widehat\pi\).  Its radial projection belongs to
\(\Piopt(\nu_\alpha,\nu_\beta)\).  The one-dimensional rearrangement
inequality therefore gives
\[
  \int r^2s^2\dd\widehat\pi
  \leq\int r^2\tau(r)^2\dd\nu_\alpha(r)=\mathsf M_2,
\]
with equality only for the monotone radial coupling.

The covariance estimate uses the weighted polar comparison proved in the
lemmas preceding \citet[Theorem~9.21]{Sturm2023Space}.  For completeness, its
key identity is
\begin{equation*}
  \ip{\omega}{\omega'}\ip{\zeta}{\zeta'}
  =\frac12\left(
    \ip{\omega}{\omega'}^2+\ip{\zeta}{\zeta'}^2
    -\abs{\ip{\omega}{\omega'}-\ip{\zeta}{\zeta'}}^2
  \right).
\end{equation*}
Combined with spherical isotropy, this identity yields the weighted
inequality
\begin{equation*}
  \iint rr'ss'\ip{\omega}{\omega'}\ip{\zeta}{\zeta'}
  \dd\widehat\pi(r,\omega,s,\zeta)
  \dd\widehat\pi(r',\omega',s',\zeta')
  \leq\frac{\mathsf M_1^2}{d}.
\end{equation*}
Its left-hand side is \(\norm{S_\pi}_{\mathrm F}^2\).  Importantly, this
weighted inequality does not follow by optimizing radii and angles
independently: an arbitrary coupling may correlate each radius with the
opposite angular variable.  Sturm's equality analysis shows that equality
forces the monotone radial coupling and
\(\ip{\omega}{\omega'}=\ip{\zeta}{\zeta'}\) for two independent copies.
The latter is zero squared GW distortion on the unit sphere, so the angular
coupling is induced by an isometry, necessarily the restriction of an
orthogonal map.  Conversely, every coupling
\eqref{eq:p2-sturm-radial-map} realizes both equalities.
\end{proof}

\begin{theorem}[Sturm's radial Brenier theorem for squared-distance GW]
\label{thm:sturm-radial-brenier}
Under the assumptions of \cref{lem:sturm-polar-maximization}, every minimizer
of squared-distance GW is induced by one of the maps \(T_O\) in
\eqref{eq:p2-sturm-radial-map}, up to changes on an \(\alpha\)-null set.
Conversely, every \(O\in\mathrm O(d)\) gives a minimizer.  Thus the optimizer
is unique modulo postcomposition by an orthogonal transformation fixing
\(m_\beta\), and
\begin{equation}
  \min_{\pi\in\Piopt(\alpha,\beta)}\Ecal_2(\pi)
  =\mathsf C_2(\alpha,\beta)
   -4\mathsf M_2-\frac{8}{d}\mathsf M_1^2.
  \label{eq:p2-sturm-radial-gw-value}
\end{equation}
\end{theorem}

\begin{proof}
For \(p=2\), \eqref{eq:gw-feature-decomposition} reads
\begin{equation*}
  \Ecal_2(\pi)
  =\mathsf C_2(\alpha,\beta)
   -4\int\norm{x-m_\alpha}^2\norm{y-m_\beta}^2\dd\pi(x,y)
   -8\norm{S_\pi}_{\mathrm F}^2.
\end{equation*}
The maps \(T_O\) simultaneously attain both upper bounds in
\eqref{eq:p2-sturm-two-bounds}, proving
\eqref{eq:p2-sturm-radial-gw-value}.  Since both coefficients are strictly
negative, any minimizer must attain both bounds.  The equality statement of
\cref{lem:sturm-polar-maximization} then gives
\(\pi=(\Id,T_O)_\sharp\alpha\), and the converse is immediate.
\end{proof}

\begin{corollary}[Isotropic Gaussian marginals]
\label{cor:p2-isotropic-gaussian-monge}
Let
\(\alpha=\mathcal N(m_\alpha,s_\alpha^2\Id)\) and
\(\beta=\mathcal N(m_\beta,s_\beta^2\Id)\) on \(\R^d\), with
\(s_\alpha,s_\beta>0\).  Every squared-distance GW optimizer is induced by
\begin{equation}
  T_O(x)
  =m_\beta+\frac{s_\beta}{s_\alpha}O(x-m_\alpha),
  \qquad O\in\mathrm O(d),
  \label{eq:p2-isotropic-gaussian-map}
\end{equation}
and
\begin{equation}
  \min_{\pi\in\Piopt(\alpha,\beta)}\Ecal_2(\pi)
  =4d(d+2)(s_\alpha^2-s_\beta^2)^2.
  \label{eq:p2-isotropic-gaussian-value}
\end{equation}
\end{corollary}

\begin{proof}
The two radial laws differ by the dilation
\(\tau(r)=(s_\beta/s_\alpha)r\), so
\cref{thm:sturm-radial-brenier} gives
\eqref{eq:p2-isotropic-gaussian-map}.  Under this coupling,
\(\norm{T_O(X)-T_O(X')}^2=(s_\beta^2/s_\alpha^2)\norm{X-X'}^2\).
Since \(X-X'\sim\mathcal N(0,2s_\alpha^2\Id)\),
\(\mathbb E[\norm{X-X'}^4]=4s_\alpha^4d(d+2)\), which proves
\eqref{eq:p2-isotropic-gaussian-value}.
\end{proof}

\subsection{The second branch is an affine reflection}

\begin{lemma}[Rank geometry of the squared-distance learned cost]
\label{lem:p2-rank-geometry}
Assume \(A\in\R^{d\times n}\) has full row rank.  For every fixed \(x\),
each fiber of
\(
 y\mapsto\nabla_xc_A(x,y)
\)
contains at most two points.  If
\begin{equation}
  \widetilde x\notin\Ran(A^\top),
  \label{eq:off-row-space}
\end{equation}
then every fiber contains at most one point, so the cost is twisted at
\(x\).

If instead \(\widetilde x\in\Ran(A^\top)\setminus\{0\}\), let
\(q_x\in\R^d\) be the unique vector satisfying
\begin{equation}
  A^\top q_x=-2\widetilde x
  \label{eq:p2-reflection-normal}
\end{equation}
and define
\begin{equation}
  \mathscr R_{A,x}(y)
  :=m_\beta+\widetilde y
    +\frac{1-2\ip{\widetilde y}{q_x}}{\norm{q_x}^2}q_x.
  \label{eq:p2-target-reflection}
\end{equation}
Then \(\mathscr R_{A,x}\) is reflection across the affine hyperplane
\begin{equation}
  m_\beta+\left\{z\in\R^d:\ip{z}{q_x}=\frac12\right\},
  \label{eq:p2-reflection-hyperplane}
\end{equation}
and, for all \(y_0,y_1\),
\begin{equation}
  \nabla_xc_A(x,y_0)=\nabla_xc_A(x,y_1)
  \quad\Longleftrightarrow\quad
  y_1\in\{y_0,\mathscr R_{A,x}(y_0)\}.
  \label{eq:p2-exact-reflection-fiber}
\end{equation}
At \(\widetilde x=0\), the gradient map is again injective.
\end{lemma}

\begin{proof}
Suppose two targets \(y_0,y\) have the same gradient.  With
\(z_0=y_0-m_\beta\), \(z=y-m_\beta\), and
\(r=\norm{z}^2-\norm{z_0}^2\), \eqref{eq:p2-cost-gradient} gives
\begin{equation}
  A^\top(z-z_0)=-2r\widetilde x.
  \label{eq:p2-gradient-fiber-equation}
\end{equation}
If \(\widetilde x\notin\Ran(A^\top)\), the two sides can agree only when
\(r=0\).  Since full row rank makes \(A^\top\) injective, this forces
\(z=z_0\).  The same conclusion follows directly when \(\widetilde x=0\).

Now suppose \(\widetilde x\in\Ran(A^\top)\setminus\{0\}\) and let \(q_x\)
be given by \eqref{eq:p2-reflection-normal}.  Equation
\eqref{eq:p2-gradient-fiber-equation} gives \(z=z_0+rq_x\).  Substitution into
the definition of \(r\) yields
\[
  r=2r\ip{z_0}{q_x}+r^2\norm{q_x}^2.
\]
Besides \(r=0\), its only possible solution is
\begin{equation*}
  r=\frac{1-2\ip{z_0}{q_x}}{\norm{q_x}^2},
\end{equation*}
which gives \(y=\mathscr R_{A,x}(y_0)\).  Conversely, for this reflected
point one has
\(
 \norm{\widetilde y}^2-\norm{\widetilde y_0}^2=r
\)
and \(A^\top(\widetilde y-\widetilde y_0)=-2r\widetilde x\), so the two
gradients agree.  Formula \eqref{eq:p2-target-reflection} is the standard
reflection formula for \eqref{eq:p2-reflection-hyperplane}; this also proves
that it is an involution and completes \eqref{eq:p2-exact-reflection-fiber}.
\end{proof}

The same full-rank algebra also gives subtwist.  This imports a uniqueness and
two-limb statement from linear OT before asking whether the second limb can be
eliminated.

\begin{corollary}[Full-rank frozen uniqueness and limb structure]
\label{cor:p2-full-rank-subtwist}
Assume compact supports and \(\alpha\ll\mathrm{Leb}^n\), and let
\(A\in\R^{d\times n}\) have full row rank.  The linear OT problem for
\(c_A\) has a unique optimizer.  This optimizer is simultaneously a two-map
and a map/anti-map.

Consequently, if \(\pi_\star\) is GW-optimal for \(p=2\) and
\(A_\star=4S_{\pi_\star}\) has full row rank, then \(\pi_\star\) has both
structures and is the unique optimizer of its frozen linear OT problem.
\end{corollary}

\begin{proof}
The two-twist assertion is \cref{lem:p2-rank-geometry}.  To verify subtwist,
fix distinct targets \(y_0,y_1\), put
\begin{equation*}
  z_i:=y_i-m_\beta,
  \qquad
  \Delta:=\norm{z_0}^2-\norm{z_1}^2,
\end{equation*}
and set \(F(x):=c_A(x,y_0)-c_A(x,y_1)\).  From
\eqref{eq:p2-learned-cost},
\begin{equation}
  \nabla F(x)=-8\Delta(x-m_\alpha)-4A^\top(z_0-z_1),
  \qquad
  \nabla^2F=-8\Delta\Id.
  \label{eq:p2-subtwist-difference}
\end{equation}
If \(\Delta=0\), injectivity of \(A^\top\) and \(z_0\ne z_1\) show that
\(F\) has no critical point.  If \(\Delta\ne0\), it has the unique critical
point
\begin{equation*}
  x_c=m_\alpha-\frac{A^\top(z_0-z_1)}{2\Delta},
\end{equation*}
which is a global maximum when \(\Delta>0\) and a global minimum when
\(\Delta<0\).  Hence \(c_A\) satisfies \cref{def:subtwist}.
\Cref{thm:subtwist-linear-ot} gives uniqueness and the map/anti-map structure,
while \cref{prop:finite-twist-linear-ot} gives the two-map structure.  The GW
claim follows from \cref{cor:self-consistent-linear-ot}.
\end{proof}

The subtwist and map/anti-map conclusions recover the full-rank part of
\citet[Theorem~3.6]{DumontLacombeVialard2025}.  The additional contribution of
\cref{lem:p2-rank-geometry} is to identify their possible second twist branch
as the explicit source-dependent reflection \eqref{eq:p2-target-reflection}.

Frozen uniqueness must not be confused with uniqueness of the nonconvex GW
problem: two GW optimizers can have different cross-covariances and therefore
different self-consistent frozen costs.  The reflection formula now gives an
exact contact-set criterion under which the unique frozen map/anti-map reduces
to a single graph.

\begin{proposition}[Full-rank reflection criterion]
\label{prop:p2-full-rank-reflection-criterion}
Let \(\pi_\star\) be GW-optimal for \(p=2\), assume compact supports and
\(\alpha\ll\mathrm{Leb}^n\), and suppose
\(A_\star=4S_{\pi_\star}\) has full row rank.  Let
\((f_\star,g_\star)\) be a dual optimizer for the frozen cost and write
\begin{equation*}
  S_x:=\bigl\{y\in\supp\beta:
  f_\star(x)+g_\star(y)=c_{A_\star}(x,y)\bigr\}.
\end{equation*}
Define the contact and support bad sets by
\begin{align}
  \mathcal B^{f_\star,g_\star}_{A_\star}
  &:=\bigl\{x\in\supp\alpha\cap(m_\alpha+\Ran(A_\star^\top)),\
       x\ne m_\alpha:\notag\\[-1mm]
  &\hspace{28mm}\text{there exists }y\in S_x\text{ such that }
       \mathscr R_{A_\star,x}(y)\in S_x\setminus\{y\}\bigr\},
       \label{eq:p2-bad-contact-reflection-set}\\
  \mathcal B_{A_\star}(\beta)
  &:=\bigl\{x\in\supp\alpha\cap(m_\alpha+\Ran(A_\star^\top)),\
       x\ne m_\alpha:\notag\\[-1mm]
  &\hspace{28mm}\text{there exists }y\in\supp\beta\text{ such that }
       \mathscr R_{A_\star,x}(y)
       \in\supp\beta\setminus\{y\}\bigr\}.
  \label{eq:p2-bad-reflection-source-set}
\end{align}
These sets are universally measurable and
\(\mathcal B^{f_\star,g_\star}_{A_\star}
 \subset\mathcal B_{A_\star}(\beta)\).  If
\(\alpha(\mathcal B^{f_\star,g_\star}_{A_\star})=0\), then \(\pi_\star\) is
induced by a map.  It is already the unique optimizer of its frozen linear OT
problem by \cref{cor:p2-full-rank-subtwist}.  In particular, the support-level condition
\(\alpha(\mathcal B_{A_\star}(\beta))=0\) is sufficient.
\end{proposition}

\begin{proof}
Because the cost and dual potentials may be chosen continuous on the compact
supports, both bad sets are projections of Borel sets and hence analytic, so
they are universally measurable.  Their inclusion is immediate from
\(S_x\subset\supp\beta\).

The optimizer \(\pi_\star\) solves linear OT for \(c_{A_\star}\) by
\cref{cor:self-consistent-linear-ot}.  Outside
\(m_\alpha+\Ran(A_\star^\top)\), and at \(x=m_\alpha\), the gradient map is
injective by \cref{lem:p2-rank-geometry}.  At every remaining source point
outside \(\mathcal B^{f_\star,g_\star}_{A_\star}\),
\eqref{eq:p2-exact-reflection-fiber} is injective on the dual contact fiber
\(S_x\).  The frozen cost is therefore twisted on its contact set at
\(\alpha\)-almost every source point.  Apply \cref{prop:twist-linear-ot}; the
unconditional frozen uniqueness is \cref{cor:p2-full-rank-subtwist}.
\end{proof}

The reflection law is particularly explicit on the line.  It shows that the
two-map theorem does not leave two arbitrary branches: their midpoint is fixed
pointwise by the optimal cross-covariance.

\begin{corollary}[One-dimensional branch symmetry]
\label{cor:p2-one-dimensional-branch-symmetry}
Let \(n=d=1\), assume compact supports and
\(\alpha\ll\mathrm{Leb}^1\), and let \(\pi_\star\) be GW-optimal for \(p=2\)
with \(S_{\pi_\star}\ne0\).  Then \(\pi_\star\) is a two-map.  Moreover, for
\(\alpha\)-almost every \(x\ne m_\alpha\), if its conditional support consists
of two distinct points \(T_1(x)\) and \(T_2(x)\), then
\begin{equation}
  \frac{T_1(x)+T_2(x)}2
  =m_\beta-\frac{S_{\pi_\star}}{x-m_\alpha}.
  \label{eq:p2-one-dimensional-branch-midpoint}
\end{equation}
\end{corollary}

\begin{proof}
Here \(A_\star=4S_{\pi_\star}\ne0\), so the two-map statement follows from
\cref{cor:p2-full-rank-subtwist}.  Equation \eqref{eq:p2-reflection-normal} reads
\(
 q_x=-(x-m_\alpha)/(2S_{\pi_\star})
\).
The reflecting hyperplane \eqref{eq:p2-reflection-hyperplane} is therefore the
single target point
\(
 m_\beta-S_{\pi_\star}/(x-m_\alpha)
\).
Distinct points in the same gradient fiber are reflections about this point,
which is exactly \eqref{eq:p2-one-dimensional-branch-midpoint}.
\end{proof}

\subsection{A global one-dimensional Monge criterion}

The reflection law describes the local shape of a frozen optimizer, but does
not by itself decide whether a single branch minimizes the original quadratic
GW functional.  On the line, the decomposition
\eqref{eq:gw-feature-decomposition} yields a checkable global criterion.  Let
\(Q_\alpha,Q_\beta:(0,1)\to\R\) be the quantile functions and set
\begin{equation}
  q_\alpha(t):=Q_\alpha(t)-m_\alpha,
  \qquad
  q_\beta(t):=Q_\beta(t)-m_\beta.
  \label{eq:p2-centered-quantiles}
\end{equation}
If \(U\) is uniform on \((0,1)\), define the increasing and decreasing
rearrangements
\begin{equation}
  \pi_+:=\operatorname{Law}\bigl(Q_\alpha(U),Q_\beta(U)\bigr),
  \qquad
  \pi_-:=\operatorname{Law}\bigl(Q_\alpha(U),Q_\beta(1-U)\bigr),
  \label{eq:p2-monotone-couplings}
\end{equation}
and their centered covariances
\begin{equation}
  S_+:=\int_0^1q_\alpha(t)q_\beta(t)\dd t,
  \qquad
  S_-:=\int_0^1q_\alpha(t)q_\beta(1-t)\dd t.
  \label{eq:p2-extreme-quantile-covariances}
\end{equation}
Two measurable functions \(r,s\) on \((0,1)\) are called
\emph{comonotone} when
\(
 (r(t)-r(u))(s(t)-s(u))\geq0
\)
for almost every \((t,u)\).  This is the sufficient equality condition used
below to attain the upper rearrangement bound.

For \(\pi\in\Piopt(\alpha,\beta)\), define the radial mixed moment and scalar
cross-covariance
\begin{equation}
  R_\pi:=\int (x-m_\alpha)^2(y-m_\beta)^2\dd\pi(x,y),
  \qquad
  S_\pi:=\int (x-m_\alpha)(y-m_\beta)\dd\pi(x,y).
  \label{eq:p2-one-dimensional-two-moments}
\end{equation}
The following exact reduction isolates the one-dimensional optimization that
remains after the frozen-cost analysis.  It is the moment-slice viewpoint
underlying the one-dimensional study of
\citet{DumontLacombeVialard2025}.

\begin{proposition}[Exact covariance-slice reduction]
\label{prop:p2-one-dimensional-scalar-reduction}
The attainable covariances form the interval
\(\mathcal I=[S_-,S_+]\).  For \(s\in\mathcal I\), define
\begin{equation}
  \mathcal R(s):=
  \max\bigl\{R_\pi:\pi\in\Piopt(\alpha,\beta),\ S_\pi=s\bigr\}.
  \label{eq:p2-covariance-slice-value}
\end{equation}
Then \(\mathcal R\) is finite, concave, and upper semicontinuous on
\(\mathcal I\).  Moreover,
\begin{equation}
  \Ecal_2(\pi)
  =\mathsf C_2(\alpha,\beta)-4R_\pi-8S_\pi^2
  \label{eq:p2-one-dimensional-objective}
\end{equation}
and hence
\begin{equation}
  \min_{\pi\in\Piopt(\alpha,\beta)}\Ecal_2(\pi)
  =\mathsf C_2(\alpha,\beta)
   -\max_{s\in[S_-,S_+]}\bigl\{4\mathcal R(s)+8s^2\bigr\}.
  \label{eq:p2-covariance-scalar-reduction}
\end{equation}
A coupling \(\pi_\star\) is GW-optimal if and only if it maximizes
\(R_\pi\) on the slice \(S_\pi=S_{\pi_\star}\) and its covariance maximizes
the scalar objective in \eqref{eq:p2-covariance-scalar-reduction}.
\end{proposition}

\begin{proof}
The coupling set is compact and convex, while \(\pi\mapsto S_\pi\) and
\(\pi\mapsto R_\pi\) are continuous and affine.  Its covariance image is
therefore a compact interval.  The rearrangement inequality identifies its
endpoints as \(S_-\) and \(S_+\) \citep{Santambrogio2015}, and compactness
gives the maximum in \eqref{eq:p2-covariance-slice-value}.

For \(s_0,s_1\in\mathcal I\), mix maximizers on the two corresponding slices.
The resulting coupling has covariance \((1-\theta)s_0+\theta s_1\) and radial
moment \((1-\theta)\mathcal R(s_0)+\theta\mathcal R(s_1)\).  Thus
\(\mathcal R\) is concave.  If \(s_k\to s\), first select a subsequence
realizing \(\limsup_k\mathcal R(s_k)\) and choose a maximizer \(\pi_k\) on
each corresponding slice.  Compactness gives a further subsequence
\(\pi_k\rightharpoonup\pi\); continuity then yields
\(S_\pi=s\) and
\(\limsup_k\mathcal R(s_k)=R_\pi\leq\mathcal R(s)\).  Hence \(\mathcal R\) is upper
semicontinuous.  Formula \eqref{eq:p2-one-dimensional-objective} is
\eqref{eq:gw-feature-decomposition} in dimension one.  Minimizing it first on
each covariance slice proves \eqref{eq:p2-covariance-scalar-reduction} and the
last assertion.
\end{proof}

\begin{theorem}[Quantile criterion for a global one-dimensional Monge optimizer]
\label{thm:p2-one-dimensional-quantile-monge}
Let \(n=d=1\), assume compact supports, and suppose
\(\alpha\ll\mathrm{Leb}^1\).
\begin{enumerate}[label=\textup{(\roman*)}]
  \item If \(q_\alpha^2\) and \(q_\beta^2\) are comonotone and
  \begin{equation}
    \abs{S_+}\geq\abs{S_-},
    \label{eq:p2-increasing-quantile-criterion}
  \end{equation}
  then \(\pi_+\) is a global minimizer of squared-distance GW.  If the
  inequality in \eqref{eq:p2-increasing-quantile-criterion} is strict, then
  \(\pi_+\) is the unique GW minimizer.

  \item If the functions \(t\mapsto q_\alpha(t)^2\) and
  \(t\mapsto q_\beta(1-t)^2\) are comonotone and
  \begin{equation}
    \abs{S_-}\geq\abs{S_+},
    \label{eq:p2-decreasing-quantile-criterion}
  \end{equation}
  then \(\pi_-\) is a global minimizer.  It is the unique GW minimizer when
  the inequality is strict.
\end{enumerate}
In either case the optimizer is induced by the corresponding increasing or
decreasing rearrangement map.
\end{theorem}

\begin{proof}
The rearrangement inequality applied first to the centered variables and
then to their squares gives, for every coupling,
\begin{equation}
  S_-\leq S_\pi\leq S_+,
  \qquad
  R_\pi\leq
  \int_0^1 Q_{(X-m_\alpha)^2}(t)
               Q_{(Y-m_\beta)^2}(t)\dd t,
  \label{eq:p2-one-dimensional-rearrangement-bounds}
\end{equation}
where \(X\sim\alpha\), \(Y\sim\beta\), and \(Q_Z\) denotes the quantile
function of a real random variable \(Z\).  In particular,
\(
 S_\pi^2\leq\max\{S_-^2,S_+^2\}
\).
The product coupling has zero covariance, so \(S_-\leq0\leq S_+\).

Under the assumptions of \textup{(i)}, the pair
\((q_\alpha(U)^2,q_\beta(U)^2)\) is comonotone, so \(\pi_+\) attains the
upper bound for \(R_\pi\).  Condition
\eqref{eq:p2-increasing-quantile-criterion} says that it also maximizes
\(S_\pi^2\).  It therefore minimizes \eqref{eq:p2-one-dimensional-objective},
or equivalently maximizes the scalar objective in
\eqref{eq:p2-covariance-scalar-reduction}.
If the covariance inequality is strict, equality in the objective forces
\(S_\pi=S_+\).  Since \(\alpha\) is atomless, equality in the upper
rearrangement inequality uniquely selects the increasing coupling \(\pi_+\).
This proves uniqueness.  The same argument, with \(U\) paired with \(1-U\),
proves \textup{(ii)}.  Finally, atomlessness of \(\alpha\) represents these
couplings by
\(
 T_+(x)=Q_\beta(F_\alpha(x))
\)
and
\(
 T_-(x)=Q_\beta(1-F_\alpha(x))
\)
for \(\alpha\)-almost every \(x\).
\end{proof}

\begin{corollary}[Symmetric one-dimensional marginals]
\label{cor:p2-one-dimensional-symmetric-monge}
Assume that \(\alpha\ll\mathrm{Leb}^1\) and \(\beta\) have compact support
and are symmetric around \(m_\alpha\) and \(m_\beta\), respectively.  Then
both \(\pi_+\) and \(\pi_-\) in \eqref{eq:p2-monotone-couplings} are
squared-distance GW minimizers and are induced by maps.
\end{corollary}

\begin{proof}
Symmetry gives
\(
 q_\alpha(1-t)=-q_\alpha(t)
\)
and
\(
 q_\beta(1-t)=-q_\beta(t)
\)
for almost every \(t\).  Since each centered quantile is increasing, both
\(q_\alpha(t)^2\) and \(q_\beta(t)^2\) are nondecreasing functions of
\(\abs{t-\tfrac12}\); they are therefore comonotone.  Moreover,
\(S_-=-S_+\).  Both parts of
\cref{thm:p2-one-dimensional-quantile-monge} apply.
\end{proof}

When \(\beta\) also has a density,
\cref{cor:p2-one-dimensional-symmetric-monge} is the \(d=1\) instance of
\cref{thm:sturm-radial-brenier}, because
\(\mathrm O(1)=\{-1,+1\}\).  The quantile criterion is more flexible on the
line: it permits a singular target and nonsymmetric pairs whenever one
orientation simultaneously optimizes the two moments.

The theorem is global: unlike self-consistent frozen optimality, its
conclusion compares the candidate with every coupling.  Its assumptions
involve only the two marginals.  For example, if
\(
 \beta=(x\mapsto m_\beta+s(x-m_\alpha))_\sharp\alpha
\)
with \(s\ne0\), then the corresponding affine map is GW-optimal.  Indeed,
the relevant squared quantiles are proportional, and Cauchy--Schwarz shows
that its covariance has maximal absolute value.  Uniqueness is governed by
the strict comparison in
\eqref{eq:p2-increasing-quantile-criterion} or
\eqref{eq:p2-decreasing-quantile-criterion}; a symmetric marginal may admit
both orientations.

These conditions are sufficient, not necessary.  In particular, the radial
comonotonicity can fail even for smooth one-dimensional marginals, leaving
room for a genuinely two-valued optimizer.  A different global sufficient
condition is proved by \citet[Proposition~3.14]{DumontLacombeVialard2025}:
for two-component measures with matching component masses and sufficiently
large separation, one of the two monotone rearrangements is the unique GW
optimizer.  Thus the quantile criterion and the separation theorem isolate
two concrete regimes where the frozen two-branch obstruction collapses to a
Monge solution, without claiming that this happens for arbitrary marginals.

\subsection{Dimension-lowering Monge theorems}

The following dimension gap makes the bad set automatically null.

The exceptional set in \eqref{eq:off-row-space} is a proper affine subspace
as soon as the source dimension is strictly larger than the target dimension.
This yields a concrete Monge theorem.

\begin{theorem}[Full-rank dimension-gap theorem for \(p=2\)]
\label{thm:p2-dimension-gap-monge}
Let \(n>d\), let \(\alpha\ll\mathrm{Leb}^n\), and assume compact supports.
Let \(\pi_\star\) be a minimizer of the squared-distance problem
\eqref{eq:power-gw-kantorovich} with \(p=2\).  If
\begin{equation}
  \rank S_{\pi_\star}=d,
  \label{eq:full-row-rank-cross-covariance}
\end{equation}
then \(\pi_\star\) is induced by a map:
\(
 \pi_\star=(\Id,T_\star)_\sharp\alpha.
\)
More precisely, \(\pi_\star\) is the unique optimizer of linear OT for
\(c_{A_\star}\), where \(A_\star=4S_{\pi_\star}\).  If \(f_\star\) is the
corresponding source potential, then
\begin{equation}
  \nabla f_\star(x)
  =-8(x-m_\alpha)\norm{T_\star(x)-m_\beta}^2
   -16S_{\pi_\star}^\top(T_\star(x)-m_\beta)
  \label{eq:p2-dimension-gap-map-characterization}
\end{equation}
for \(\alpha\)-almost every \(x\).
\end{theorem}

\begin{proof}
By \cref{cor:self-consistent-linear-ot}, \(\pi_\star\) solves linear OT for
the cost \eqref{eq:p2-learned-cost} with
\(A_\star=4S_{\pi_\star}\).  Assumption
\eqref{eq:full-row-rank-cross-covariance} makes \(A_\star\) full row rank.
The set \(\mathcal B_{A_\star}(\beta)\) in
\eqref{eq:p2-bad-reflection-source-set} is contained in the
\(d\)-dimensional affine space \(m_\alpha+\Ran(A_\star^\top)\).  Since
\(d<n\),
\[
  \alpha\bigl(\mathcal B_{A_\star}(\beta)\bigr)=0.
\]
Apply \cref{prop:p2-full-rank-reflection-criterion}.  Substituting
\(A_\star=4S_{\pi_\star}\) into \eqref{eq:p2-cost-gradient} gives
\eqref{eq:p2-dimension-gap-map-characterization}.
\end{proof}

This conclusion is stronger than the generic count in
\cref{prop:general-position-branch-count}.  When \(d<n\leq2d\), that count
allows several branches, whereas the special rank-one quadratic structure of
\eqref{eq:p2-cost-gradient} confines all branching to the \(d\)-dimensional
affine space \(m_\alpha+\Ran(A_\star^\top)\), which is \(\alpha\)-null.

The dimension gap is not the only way to remove the second branch.  If the
target feature radius is constant, the nonlinear radial term in
\eqref{eq:p2-learned-cost} becomes marginal-only and the learned cost is
effectively bilinear.

\begin{corollary}[Centered spherical targets]
\label{cor:p2-spherical-target}
Let \(n\geq d\), assume \(\alpha\ll\mathrm{Leb}^n\), and let
\(\pi_\star\) be a minimizer for \(p=2\) with
\(\rank S_{\pi_\star}=d\).  If, for some \(R\geq0\),
\begin{equation}
  \norm{y-m_\beta}=R
  \qquad\text{for every }y\in\supp\beta,
  \label{eq:centered-spherical-target}
\end{equation}
then \(\pi_\star=(\Id,T_\star)_\sharp\alpha\), including when \(n=d\).
\end{corollary}

\begin{proof}
Suppose \(y\) and \(\mathscr R_{A_\star,x}(y)\) both belong to
\(\supp\beta\).  Their centered norms are equal by
\eqref{eq:centered-spherical-target}.  In the proof of
\cref{lem:p2-rank-geometry}, the reflection coefficient is exactly the
difference of their squared centered norms, so it vanishes and the two points
coincide.  Thus \(\mathcal B_{A_\star}(\beta)=\varnothing\), and
\cref{prop:p2-full-rank-reflection-criterion} applies.
\end{proof}

\begin{remark}[What the theorem does and does not assume]
The rank conditions in \cref{thm:p2-dimension-gap-monge,cor:p2-spherical-target}
are imposed on an optimal cross-covariance, hence they are not a priori
conditions on the marginals alone.  Their conclusions nevertheless concern
the original GW optimizer, not a Gaussian or affine restriction.  Under the
strict dimension gap of \cref{thm:p2-dimension-gap-monge}, if every optimal
cross-covariance has full row rank, then every GW optimizer is graph-supported;
the same statement holds under the spherical-target hypothesis of
\cref{cor:p2-spherical-target}.  Full row rank alone is not claimed to suffice
in equal dimensions for a general target.
\end{remark}

The preceding direct-twist theorem and the quotient construction cover
complementary rank regimes.  Their combination removes the rank hypothesis
altogether when the source dimension is strictly larger than the target
dimension.

\begin{theorem}[Dimension-lowering Brenier theorem for squared-distance GW]
\label{thm:p2-strict-dimension-gap-monge}
Let \(n>d\), let \(\alpha\in\Pcal(\R^n)\) and
\(\beta\in\Pcal(\R^d)\) have compact support, and assume
\(\alpha\ll\mathrm{Leb}^n\).  Then the squared-distance problem
\eqref{eq:power-gw-kantorovich} with \(p=2\) admits a Monge minimizer.  In
particular,
\begin{equation}
  \GW_{2,2}^2(\alpha,\beta)
  =\GM_{2,2}^2(\alpha,\beta)
  =\min_{T_\sharp\alpha=\beta}
    \Ecal_2\bigl((\Id,T)_\sharp\alpha\bigr).
  \label{eq:p2-dimension-gap-gw-gm-equality}
\end{equation}
More precisely, fix any GW minimizer \(\pi_\star\) and put
\(h=\rank S_{\pi_\star}\).
\begin{enumerate}[label=(\roman*)]
  \item If \(h=d\), then \(\pi_\star\) itself is induced by a map and is the
  unique optimizer of its frozen linear OT problem.
  \item If \(h<d\), there exists a GW-optimal map \(T_\star\) such that
  \begin{equation}
    S_{(\Id,T_\star)_\sharp\alpha}=S_{\pi_\star}.
    \label{eq:p2-dimension-gap-preserved-cross-covariance}
  \end{equation}
\end{enumerate}
\end{theorem}

\begin{proof}
If \(h=d\), \cref{thm:p2-dimension-gap-monge} proves the first assertion.
If \(h<d\), then \(h\leq d-1\) and \(n>d\geq h+1\), so in fact
\(n>h+1\).  Corollary~\ref{cor:p2-finite-rank-quotient} therefore provides a
GW-optimal graph coupling with the same cross-covariance as \(\pi_\star\).
The two cases exhaust all possible ranks and prove existence of a Monge
minimizer.  The last equality in
\eqref{eq:p2-dimension-gap-gw-gm-equality} follows because this graph coupling
attains the Kantorovich minimum, while the reverse inequality is immediate
from the inclusion of graph couplings among all couplings.
\end{proof}

The lower-rank ingredient in this proof is the quotient construction of
\citet{DumontLacombeVialard2025}.  The full-rank twist argument and the
rank-exhaustion step above make explicit how that result closes with the
complementary rank regime under the strict dimension gap.

This theorem is directional.  If \(d>n\) and instead
\(\beta\ll\mathrm{Leb}^d\), exchanging the two spaces produces a GW-optimal
map from \(\beta\) to \(\alpha\).  Its content is also stronger than equality
of the two infimum values: it proves that the Gromov--Monge infimum is
attained.  The word ``Brenier'' is used here in this structural sense.  In the
lower-rank case, the quotient map \(t_\star\) is a genuine Brenier map after
the invertible linear change \(D_{A_\star}\), but the measurable transport
inside its fibers is not canonical; consequently the lifted ambient map
\(T_\star\) need not be the gradient of a convex function on \(\R^n\).

The proof clarifies the role of the extra source dimensions.  For
\(h<d\), they provide positive-dimensional, nonatomic level sets of the
augmented statistic \(\Phi_{A_\star}\) in
\eqref{eq:p2-source-quotient-map},
which supply the randomness needed by the deterministic lift.  For \(h=d\),
the exceptional set on which the frozen cost has a second twist branch is the
\(d\)-dimensional affine space
\(m_\alpha+\Ran(S_{\pi_\star}^\top)\), hence is \(\alpha\)-null.

At the equal-dimensional boundary \(n=d\), the same proof yields a Monge
optimizer when \(h\leq d-2\), recovering the corresponding result of
\citet{DumontLacombeVialard2025}.  For \(h=d-1\), the quotient fibers become
zero dimensional; for \(h=d\), the exceptional row space fills the source.
The general guarantee then falls back to the published two-map theorem.  A
centered spherical target removes the radial branch and restores a map in the
full-rank case by \cref{cor:p2-spherical-target}; more generally,
\cref{prop:p2-full-rank-reflection-criterion} reduces this case to the absence
of source-dependent reflected pairs on the dual contact fibers.  Thus strict
dimension reduction is exactly what closes every rank case in the argument
above; it does not contradict the one-dimensional numerical evidence for
genuinely two-valued optimizers.

\subsection{Why equal dimensions retain two branches}

When \(n=d\) and \(A\) is invertible, the second part of
\cref{lem:p2-rank-geometry} applies away from the single point
\(x=m_\alpha\).  The second inverse branch is not merely an unspecified root
of a scalar quadratic: it is exactly the affine reflection
\(\mathscr R_{A,x}\) in \eqref{eq:p2-target-reflection}.  Complementary
slackness therefore restricts each conditional plan to a point and its
reflection, but does not force the latter to be absent.  The sharp contact-set
criterion is \cref{prop:p2-full-rank-reflection-criterion}; it is the local
mechanism behind \cref{thm:known-two-map}.  At rank \(d-1\),
\cref{prop:p2-injective-target-feature} identifies a different but related
obstruction: the second branch survives only when the target support contains
a pair with the same radial-linear feature, namely a pair reflected across the
fixed affine space \(m_\beta+E_\star\).

The distinction with inner-product GW is now exact.  Its learned cost is
bilinear, so the gradient equation is linear in \(y\), and full rank gives a
single inverse.  Squared-distance GW adds the radial statistic
\(\norm{y-m_\beta}^2\); solving simultaneously for the linear feature and
that radial statistic produces a quadratic scalar equation.  An RKHS norm
penalty controls the learned cost but cannot remove this algebraic second
branch.

\section{What remains for \texorpdfstring{\(0<p<2\)}{0<p<2}}
\label{sec:open}

The learned-cost theorem is exact for every \(0<p\leq2\), but the sharp
two-branch algebra used at \(p=2\) does not survive unchanged.  The preceding
finite-twist and quotient-lifting frameworks narrow the missing step without
pretending that regularity alone solves it.

\paragraph{The feature lift is not a free Brenier theorem.}
For \(0<p<2\), the Schoenberg feature space is generally infinite
dimensional.  Even when \(\alpha\) has a smooth density in \(\R^n\), its
feature pushforward \(\Phi_\sharp\alpha\) lies on a nonlinear
finite-dimensional subset of \(\Hcal_{p,n}\).  It is not absolutely
continuous with respect to an infinite-dimensional analogue of Lebesgue
measure.  Applying a Hilbert-space version of Brenier's theorem directly to
the feature measures is therefore unjustified.  The finite-rank quotient
theorem avoids this mistake by retaining only \(R+1\) visible statistics and
by explicitly transporting within the discarded fibers.  It does not apply
automatically when the optimal operator has infinite rank.

\paragraph{There are two targets: finite twist or a liftable quotient.}
For a given optimizer, \eqref{eq:power-distance-twist-map} is the explicit
twist object to study.  Uniformly bounding its contact fibers yields a
finite-map optimizer by \cref{cor:distance-only-twist}; injectivity is only
the special case that yields a Monge map.  Alternatively, a
finite-dimensional factorization of the frozen cost can produce a Monge
optimizer through \cref{thm:quotient-fiber-gw}, even when the full twist
fibers are positive dimensional.  The cross-Hessian condition
\eqref{eq:power-cross-hessian-rank-condition} proves such finiteness by local
immersion and compactness.  More flexibly,
\cref{prop:general-position-branch-count} gives the candidate generic count
\(m_{n,d}\) when \(n>d\).  What remains is to derive either finite twist or
a finite-dimensional liftable quotient from assumptions on the marginals
while respecting self-consistency of \(\pi_\star\).

\paragraph{Equal dimensions require rigidity beyond dimension counting.}
When \(n=d\), multijet dimension counting gives no finite universal branch
number.  Polynomial degree controls the \(p=2\) model, but for
\(0<p<2\) the frozen cost is generally neither polynomial nor analytic on the
diagonal.  The definable criterion of
\cref{prop:definable-finite-twist} becomes effective only after one proves
that positive-dimensional fibers are absent.  Analyticity by itself cannot
provide a dimension-only value of \(m\), as the Chebyshev example shows.

\paragraph{The nonsmooth range needs a subdifferential version.}
For \(1<p<2\), the mixed derivative in
\eqref{eq:power-mixed-hessian-identity} remains continuous and the finite
twist program is directly available.  At \(p\leq1\), the power distance is
nondifferentiable or singular at coincidences.  One must replace
\(\nabla_xc\) by a suitable subdifferential fiber and control the resulting
set-valued contact relation; this note does not yet prove that extension.

\paragraph{The endpoint \(p>2\) requires another mechanism.}
Euclidean \(p\)-powers cease to be conditionally negative definite in
general when \(p>2\) \citep{Schoenberg1938}.  The Hilbert feature
representation, the Hilbert--Schmidt cross-covariance, and the quadratic RKHS
regularizer used here then disappear.  A Monge theorem in that range would
need a different factorization or a direct analysis of the frozen GW cost.

\subsection*{Conclusion}

The strongest rigorous narrative is therefore not that the RKHS reduction
automatically proves a universal Brenier theorem.  It proves something more
precise:
\begin{enumerate}[label=\arabic*.]
  \item for a nonatomic source, graph couplings are dense and the directed
  Gromov--Monge and GW values agree, but the deterministic infimum need not be
  attained;
  \item for \(0<p\leq2\), power-distance GW is exactly regularized learning
  of a cross-space cost in a tensor-product RKHS;
  \item every GW optimizer solves ordinary OT for its self-consistent learned
  cost;
  \item for every quadratic GW kernel, frozen optimality is a universal
  first-order condition; conditional concavity upgrades it to a replacement
  principle for all frozen optimizers;
  \item an \(m\)-twisted contact relation forces a prescribed optimizer to be a
  measurable mixture of at most \(m\) maps, while one-twist gives a unique
  frozen graph coupling;
  \item subtwist gives a complementary uniqueness mechanism: every full-rank
  squared-distance frozen problem has a unique optimizer that is simultaneously
  a two-map and a map/anti-map, although the GW problem itself need not be
  unique;
  \item a twisted quotient problem together with nonatomic source fibers
  yields a GW-optimal graph coupling even when the full cost has infinite
  twist fibers;
  \item for a rank-\(R\) learned operator, this quotient has only \(R+1\)
  coordinates, and the coarea criterion \(n>R+1\) makes its liftability
  explicit;
  \item full column rank of the power cross-Hessian gives some finite branch
  number, and fiberwise general position predicts the explicit count
  \(m_{n,d}=\lfloor n/(n-d)\rfloor\) when \(n>d\);
  \item for \(p=2\), quotient lifting handles every rank below \(d\), while
  the special rank-one structure of the full-rank frozen twist gives one
  branch almost everywhere when \(n>d\); consequently a strict source-target
  dimension gap guarantees a GW-optimal map without any rank assumption;
  \item under full row rank, the possible second target is exactly an affine
  reflection depending on the source point; excluding reflected pairs on the
  dual contact fibers is an exact twist-based map criterion, and a centered
  spherical target satisfies it automatically;
  \item for equal-dimensional rotationally invariant absolutely continuous
  marginals, Sturm's theorem completely resolves the global problem: every
  optimizer is the monotone radial quantile map followed by an orthogonal
  transformation;
  \item in one dimension, centered quantiles give a more flexible
  marginal-only criterion, allowing singular targets and nonsymmetric pairs:
  if one monotone rearrangement simultaneously maximizes the radial mixed
  moment and covariance magnitude, it is GW-optimal, and a strict covariance
  comparison makes it unique;
  \item at lower rank, injectivity of the radial-linear target feature gives a
  complementary map criterion, notably at rank \(d-1\) in equal dimensions;
  and
  \item otherwise, in equal dimensions the reflection branch can survive,
  consistently with the known two-map theorem.
\end{enumerate}
The remaining research problem for \(1<p<2\) is therefore concrete: prove
nondegeneracy, transversality, finite definable fibers, or a finite-rank
liftable quotient for the self-consistent field \(G_{\pi_\star,x}\).  A
one-map theorem through twist requires injectivity; quotient lifting instead
requires sufficient conditional mass inside the invisible fibers.

\begingroup
\small
\setlength{\bibsep}{2pt plus 1pt}
\bibliographystyle{plainnat}
\bibliography{references}
\endgroup

\end{document}
